\subsection{八阶精确系数}
\label{eighth:section}

八阶价格同时取决于目标缺形的第二次反馈、三近面源层和实际删除质量。
我们先保留一个定量粗上包，由竞争输入提取正预算稳定性；
随后在整个有效源盒内计算联合高度，最后用同一八层账本消去源缺量。
二维的微孔还产生独立的小半径项，须在真实基线导数下单独支付。

\begin{theorem}[八阶精确展开]\label{eighth:theorem}
固定整数 \(n\ge2\)，令 \(a=v^n=2^nt\)、\(v=2t^{1/n}\)，并记
\[
 K_n=\binom n2\frac{n^2}{27},\qquad
 P_{7,n}(t)=L_n(t)+K_nav^6-\left(n+\frac{26}{9}\right)K_nav^7,
\]
\[
 J_n=\frac{(n-1)n^3(2609+1971n+81n^2)}{26244}
       +\frac8{27}\mathbf1_{\{n=2\}},\qquad
 \widehat J_n=J_n+\frac8{27}\mathbf1_{\{n=2\}}.
\]
则
\begin{equation}\label{eighth:exact-expansion}
 C_n(1-t)=P_{7,n}(t)+\widehat J_nav^8+o_n(av^8).
\end{equation}
特别 \(\widehat J_2=17638/6561\)。所有余项按固定维数理解。
\end{theorem}

\subsubsection{同一八层族与定量粗先验}

记 \(A_v=(1+v)^n\)、\(B_n=\binom n2\)、\(k=n^2/27\)、
\(\Gamma=(10/9)k\)。上界只须考虑任意序列
\begin{equation}\label{eighth:critical}
 \begin{aligned}
 &v\downarrow0,\quad V=|E|>0,\quad\\
 &K\subset\mathcal H_n(1-t;E)\setminus E\ \text{紧},\quad M=|K|,\\
 &\qquad 1+M/V\ge L_n(t)+K_nav^6-B_0av^7,
 \end{aligned}
\end{equation}
其中 \(B_0=(n+26/9)K_n+1\) 固定。
沿用原有限贪心族及源、目标记号，写 \(S=\sum_i r_i^n\)、\(\mathcal V=V+M\)。
所有 \(O\) 常数只依赖固定维数及固定选择参数，不依赖输入、指标数或尺度跨度；
little-o 沿任意上述输入序列成立，模量可以依赖该序列。

\begin{lemma}[同一八层族的三个前缀]\label{eighth:priors}
可先按原 \(P,S\) 图固定一个 \(I\)，使
\eqref{sharp-graph-absorption} 成立，再一次构造八层源、目标。
统一全部深度常数为 \(A_8\) 后才选择强 CDF 斜率。
对本次族的 \(m=6,7,8\)，有
\begin{align}
 P_I+\mathcal D_m/a+v\sum_i(\varepsilon_{i,m}+\mathfrak q_{i,m})r_i^n
 +\mathcal F_m+c_\tau\mathcal F_{\tau,m}&=O(v^6S),
 \label{eighth:positive-prior}\\
 \mathcal D_{m-1}/a+\mathfrak d_m^0/a+S_I+S(\partial I)&=O(v^5S),
 \qquad \Delta_m^{\rm mass}/a=O(v^6S).
 \label{eighth:mass-prior}
\end{align}
相应宽度、固定裕量稳健面及全尺度局部匹配分别在各实际末层成立。
\end{lemma}
\begin{proof}
纯维数价格取 \(c_{\rm rec}=1/(2^{n+1})\)、
\(C_{\rm b}=2C_2(n)+2\)、\(C_*=C_{\rm b}+2\)。
原图选层只用原 \(P,S\) 和这些价格，故在源深度之前固定。
引理~\ref{lem:global-exterior-union} 的有限递归执行至八层；
每步先固定深目标，再固定中心邻域，边界仍用原私有源。
同一 \(I,K',M'\) 和原指定见证用于全部层，完整本目标与异源目标排他性不变。
有精确质量恒等式 \eqref{sharp-mass-identity}，以及
\begin{equation}\label{eighth:depth}
 \mathcal D_h\ge f_h\ge0,\quad
 \mathfrak d_h^0\le f_{h-1}+aP_\partial,\quad
 f_h=M'-\sum_i|N_{i,h}|,\quad f_0=A_0vM',\qquad1\le h\le8.
\end{equation}
第一层使用 \(\mathfrak d_1^0\le f_0\)。
对 \eqref{sharp-coarse-retained-recursion} 分别取长度 \(m=6,7,8\) 的前缀，
以 \((b_0v/c_0)^{m-h}\) 加权；中间缺量精确消去，初端除 \(a\) 为
\(O(v^{m+1}S)\)。末层正量和
\(\mathcal D_m/a,v\mathcal D_{m-1}/a\) 的固定份额仍保留。
输入假设和已证粗结算给
\(\Xi=\mathcal V-A_vM/a=O(v^6S)\) 的双侧界，
同一图条件遂给所述先验。
引理~\ref{sharp-critical} 的宽度、面对和
引理~\ref{sharp-input-tiling} 的原输入证明，只用这些固定末层条件，
所以在本次三个前缀重新应用，不借用旧层对象。
特别
\begin{equation}\label{eighth:leading-mass}
 M/V=a[1+O(v)],\qquad S/\mathcal V=1+O(v),\qquad
 M/a\le S\le V/(1-a).
\end{equation}
\end{proof}

以下局部论证固定取本次第七或第八层，略去层号。
各层分别重取好根、正常槽及实际源层，有限多个常数可预先统一。
记
\[
 R_i=(g_i/a)^{1/n},\quad
 X_i=x_i+W_i+[-R_i/2,R_i/2]^n,\quad
 J_i=(R_i+|W_i|^{1/n})^n-(R_i+\mu_i^{1/n})^n,
\]
\[
 H=\sum_{i<j}|X_i\cap X_j|-\sum_iJ_i,\quad
 \mathscr D_i=(x_i+W_i^0)\setminus F_i,\quad \Delta_i=\Delta_i^{\rm end}.
\]
删除质量 \(\Delta_h^{\rm mass}\) 与这里的 CDF 端移分开。
由原强斜率选择，对任意 \(Y_i\subset(x_i+W_i')\setminus F_i\)，
\begin{equation}\label{eighth:cdf-reserve}
 |W_i|-\mu_i\ge |Y_i|+
 c\bigl[|\mathscr D_i\setminus Y_i|+(vR_i)^{n-1}\Delta_i\bigr].
\end{equation}
确实，中点与有效盒的体积差扣除源尾后仍至少为
\(c(vR_i)^{n-1}\Delta_i\)，原端盒外环至多是其固定倍；
先取出 \(Y_i\) 再分配剩余，其系数仍为一。
完整求导区间上
\begin{equation}\label{eighth:derivative}
 \frac{d}{dm}(R_i+m^{1/n})^n
 \ge(1-C\varepsilon_i)(1+v)^{n-1}\frac va,\qquad
 \frac{d}{dm}(R_i+m^{1/n})^n\ge c\,v/a.
\end{equation}
还记三个端盒的相对宽度误差为
\[
 \xi_i=\max_{l,\ W_i^0,W_i',W_i}|w_{il}/(vR_i)-1|,\qquad
 \sum_i\xi_iR_i^n=O(v^{5/2}S),\quad 0\le\xi_i\le C.
\]
这由宽度平方界、Cauchy--Schwarz 及端移先验得到。
只需上述粗界
\[
 \begin{aligned}
 &\sum\varepsilon_iR_i^n=O(v^5S)\text{、}\\
 &\sum_{\rm deep}(\alpha_i+\Delta_i/R_i)R_i^n=O(v^6S)\text{，}
 \end{aligned}
\]
不预先假定七阶饱和或定性缺源轮廓。

\subsubsection{实际覆盖、局部根传输与软接触容量}

\begin{lemma}[实际空隙及窗口根差]\label{eighth:window}
令 \(U=\bigcup_{i\notin I}X_i\)、\(\mathcal Q=\bigcup_iQ_i\)。
则
\begin{equation}\label{eighth:void}
 |\mathcal Q\setminus U|=O((v^{n+1}+v^5)S),\qquad
 \sum_{i<j}|X_i\cap X_j|=O(v^6S).
\end{equation}
对好根 \(G_v=Z_v^c\) 和任意固定 \(R\)，完整窗口清单
\(\mathcal I_i(R)=\{j:B_j\cap Q(x_i,Rr_i)\ne\varnothing\}\) 满足
\begin{equation}\label{eighth:root-transfer}
 \sum_{i\in G_v}\sum_{j\in\mathcal I_i(R)}
       |R_i-R_j|r_i^{n-1}=O_R(v^6S).
\end{equation}
\end{lemma}
\begin{proof}
第二个粗交量界由引理~\ref{sharp-angle-integrability} 的原 Hölder 估计给出。
单盒基线、质量恒等式及 Bonferroni 给
\[
 |U|\ge A_vM'/a-O(v^6S)\ge V+M-O(v^5S).
\]
原严格见证给 \(|\mathcal Q|\le V+aS\)，又
\[
 0\le S-M/a\le
 V\left[(1-a)^{-1}-(A_v-a)^{-1}+O(v^6)\right]=O(vS).
\]
因而空隙至多 \(aS-M+O(v^5S)\)，即第一式。

为证明窗口根差，先在任意单位立方铺砌中连接固定窗口内的两胞。
在各胞中心边长 \(1/4\) 的内盒独立均匀取端点，连接线段留在固定扩张窗口。
窗口所遇胞数有统一界 \(N_R\)。条件化两端点法向坐标后，
线段与某面平面的交点，其切向密度至多 \(8^{n-1}\)：
两个均匀密度卷积的界为
\(\min(4^{n-1}/(1-\theta)^{n-1},4^{n-1}/\theta^{n-1})\)。
窗口内至多 \(N_R^2\) 个面片，故可选固定
\(\tau_R=(4\cdot8^{n-1}N_R^2)^{-1}\)，使存在一条线段只穿过
面积至少 \(\tau_R\) 的面片，并避开余维至少二的边界。
每胞凸，线段不重入，路径长度有界；面片每个切向边亦至少 \(\tau_R\)。

引理~\ref{sharp-input-tiling} 的完整清单将这条路径拉回原族。
只须匹配误差小于一个固定阈值，便得到固定裕量的实际稳健面路径。
沿路精确相加实际根差，再用 \eqref{sharp-face-tightness}。
一个实际面边能出现的根均位于它的固定倍尺度窗口，尺度可比；
原半盒 packing 使重复次数有统一界。这证明 \eqref{eighth:root-transfer}，
没有把模型位移作为加法面积误差。
\end{proof}

对每个外部二临界正交边 \(e=\{i,j\}\)，令
\[
 R_e=\max(R_i,R_j),\quad
 p_e=L_1/(v^3R_e),\quad q_e=L_2/(v^3R_e),\quad m_e=\min(p_e,q_e),
\]
其余轴实际交投影记为 \(T_e\)、\(b_e=|T_e|\)。
固定 \(C\ge1\)，置 \(\omega_e=\min(1,Cm_e)\)、
\(W_C=\sum_e b_eR_e^2\omega_e\)；二维空投影的体积取一。

\begin{lemma}[自适应软权容量]\label{eighth:soft-capacity}
对任意固定 \(C\ge1\)，
\begin{equation}\label{eighth:soft-weight}
 W_C\le\frac{B_n}{(1+v)^2}\mathcal V+O_C(v^3S).
\end{equation}
\end{lemma}
\begin{proof}
先取 \(n\ge3\)，每条好端点边以较小编号端为锚，反射临界轴为右上布局，
以实际 \(X_i\) 中心为平面原点、\(r_i\) 为尺度。固定很小
\(\epsilon>0\)，放置
\[
 P_B=[1/2+\epsilon,1/2+2\epsilon]\times[1/2-2\epsilon,1/2-\epsilon],
 \qquad P_C=P_B^{\rm transpose}.
\]
每根在所有轴对、符号下只有 \(8B_n\) 个 probe。
若某 fibre 上 probe 至少半面积未被 \(U\) 覆盖，Fubini 与原 \(Q_i\)
有界重数给该类 fibre 的 \(r_i^2\) 加权量至多
\(C|\mathcal Q\setminus U|\)；所有 probe 均位于 \(Q_i\) 平面投影内。

令 \(\ell_e=\min L_l\)、\(\lambda_e=\ell_e/(1000r_i)>0\)。
原临界坐标界使所有 \(\lambda_e<\epsilon^2\)，不设共同正下界。
再删除 fibre 上清单中某对不同胞的平面交面积至少
\(\lambda_e r_i^2\) 的部分。剩余 \(\Gamma_e\subset T_e\) 称为清洁片。
清单固定取 \(\mathcal I_i(20)\)；碰 probe 的覆盖胞必与 \(Q_i\) 相交，
好根的完整窗口和图深度保证它在清单内、深处且尺度趋同。
这些均是有限盒的 Borel 判据。

清洁 fibre 上 probe 被覆盖者必是正确的实际辅助胞。
具体令归一化中心、边长误差 \(\delta\le\epsilon/100\)。
覆盖 \(P_B\) 的胞中心满足
\(\epsilon-\delta\le b_x\le1+2\epsilon+\delta\)、
\(-2\epsilon-\delta\le b_y\le1-\epsilon+\delta\)。
若 \(b_x\le1-4\epsilon\)，它与左下端胞的交面积至少
\((4\epsilon-\delta)(\epsilon-2\delta)>2\epsilon^2\)；
若 \(b_y\ge4\epsilon\)，它与右上端胞的交面积至少
\((\epsilon-3\delta)(4\epsilon-2\delta)>2\epsilon^2\)。
两者均触发已删碰撞。因此辅助胞位于 \((1,0)\) 的 \(4\epsilon\) 邻域，
另一 probe 给 \((0,1)\) 邻域的辅助胞。
同列两个不同胞的交面积大于固定正比例，也触发碰撞，故二者唯一。
按这两个有限清单指标，\(\Gamma_e\) 分成有限 Borel 片。

每片同一 transverse 点分配给四胞 \(A,B,C,D\) 的四个相应角面位置。
这些都是实际角面，四胞都包含该点。
须比较不同锚点的分配：若共享同一胞、轴对、角符号和该点，
两组其余三个中心落入同三个近单位列。经共享胞比较尺度，
第二组各胞中心到第一锚点至多 \((6+o(1))r_i\)，
原半盒半径为 \((1/2+o(1))r_i\)，故全在第一清单。
同列指标若不同，有固定大面积碰撞；所以两组实际四胞完全相同。
两条不同边只能是这四胞的相反对角。实际端点给精确恒等式
\[
 L_{AD,x}+L_{BC,x}=(A_x^+-B_x^-)+(C_x^+-D_x^-).
\]
第一边的 \(x\) 交长至少 \(\ell_e\)，第二边只须为正；
故至少一个相邻面侵入量大于 \(\ell_e/2\)。
另一平面轴交长大于 \(r_i/2\)，产生面积大于
\(\ell_er_i/4>\lambda_er_i^2\) 的第一清单碰撞，矛盾。
每个实际角面点至多用一次，此结论不要求两边共享阈值。

收费时必须先抵消每边小阈值：
\begin{equation}\label{eighth:threshold-cancel}
 \frac{\omega_e}{\lambda_e}
 \le\frac{C\ell_e/(v^3R_e)}{\ell_e/(1000r_i)}
 \le C_nC/v^3.
\end{equation}
打开碰撞并用求和上界，对清单 pair \(s,t\) 的平面交面积
\(A_{st}\ge\lambda_er_i^2\)，轴平行性给
\[
 R_e^2\omega_e|\pi_{\widehat{qr}}(X_s\cap X_t)|
 =\frac{R_e^2\omega_e}{A_{st}}|X_s\cap X_t|
 \le C_nC v^{-3}|X_s\cap X_t|.
\]
每锚点边数有界，一个全局 pair 只进入有界个可比尺度窗口。
由 \eqref{eighth:void}，总碰撞费为 \(O_C(v^3S)\)。
probe 空隙费为 \(O((v^{n+1}+v^5)S)\)，坏端点费为 \(O(v^{9/2}S)\)。
把每个分配的 \(R_e^2\) 换为所属胞 \(R_c^2\)，
误差由 \eqref{eighth:root-transfer} 支付为 \(O(v^6S)\)；
同一 helper 的片先求并，不随分片数重复收费。
角面单射和 \(\omega_e\le1\) 遂给
\[
 4W_C\le\sum_{i,\,q<r,\,\pm,\pm}
 R_i^2\prod_{l\ne q,r}(R_i+w_{il})+O_C(v^3S).
\]
宽度 \(L^1\) 误差由 \eqref{sharp-width-tightness} 为 \(O(v^{7/2}S)\)。
逐乘积展开及精确质量式给每个轴对
\[
 \begin{aligned}
 \sum_iR_i^2\prod_{l\ne q,r}(R_i+w_{il})
 &=(1+v)^{n-2}\sum_iR_i^n+O(v^{7/2}S),\qquad\\
 \sum_iR_i^n&=(M'+\mathcal D)/a\le\mathcal V/A_v+O(v^6S).
 \end{aligned}
\]
四个符号消去系数四，得到结论。

二维没有 transverse 变量，另作离散计数。
好边 \(i,j\) 的局部铺砌为中心 \((0,0),(1,1)\) 的单位方体。
覆盖 \((1/2+\eta,1/2-\eta)\) 的胞，与两指定胞分离，
依次迫使其中心横坐标至少一、纵坐标至多零；
令 \(\eta\downarrow0\)，局部有限性给唯一辅助胞 \((1,0)\)，另一胞为 \((0,1)\)。
完整匹配拉回四个实际胞，其四个相邻面有固定稳健裕量及精确 CDF 面序。
每胞每角至多分配一次：两组共享角由中心分离先迫使四胞相同，
两相反对角同时正交又被这四条精确面序排除。
将每边 \(R_e^2/4\) 分到四角，根权差用至多两条稳健面路径支付。
故无权角和至多 \(\sum_iR_i^2+O(v^{9/2}S)\le\mathcal V/A_v+O(v^{9/2}S)\)。
由 \(W_C\) 不大于无权和，即得二维情形。
\end{proof}

\subsubsection{正常槽空缺与全部高支撑的净预算}

对深处二临界边沿用 \(p,q,R_e,b_e\)，并记
\[
 \begin{aligned}
 z&=\frac{\operatorname{clip}((R_i-R_j)/2;-\max L_l,\max L_l)}{v^3R_e},\\
 &\quad x=\frac{1+v}{n}p,\quad y=\frac{1+v}{n}q,\quad
 \zeta=\frac{1+v}{n}z,\quad\\
 &\widehat b=\frac{b_e}{(1+v)^{n-2}R_e^{n-2}},\quad b=\min(1,\widehat b).
 \end{aligned}
\]
固定小 \(\nu>0\)，正常集合 \(\mathcal N\) 定义为
\begin{equation}\label{eighth:normal}
 |x-1/3|<\nu,\quad |y-1/3|<\nu,\quad|\zeta|<\nu,\quad1-b<\nu.
\end{equation}
取 \(\nu\) 足够小，正常根差未截断，\(R_i/R_j=1+O(v^3)\)，且
\[
 c_nv^3R_e\le L_1,L_2\le T_0v^3R_e,\qquad
 b_e\ge(1-\nu)R_e^{n-2}.
\]
正常两端非临界中心差小于 \(r_i/100\)。
同源同轴对同符号角至多一个正常目标，否则其全部中心坐标过近而违反原分离。

取前述同一清洁分配，令 \(O_s\subset F_s\) 为送入实际槽 \(s\) 的正常片之并。
这里槽 \(s=(i,q,r,\sigma_q,\sigma_r)\) 的横向容量为
\(F_s=\pi_{\widehat{qr}}X_i\)。
记完整槽容量和真实非负空缺为
\begin{equation}\label{eighth:slot-deficit}
 \mathcal C=\frac14\sum_sR_i^2|F_s|,\qquad
 D_{\mathcal N}=\frac14\sum_sR_i^2|F_s\setminus O_s|.
\end{equation}
正常软权为一时，漏片和根权传输给
\begin{equation}\label{eighth:slot-identities}
 W_{\mathcal N}=\mathcal C-D_{\mathcal N}+O(v^3S),\qquad
 W_{\mathcal N}+W_O\le\mathcal C+O(v^3S),
\end{equation}
其中 \(W_O\) 是其余全部二轴边的软权和，含边界。
占用先在同槽求并，不假设一条边占满 helper 整槽。

\begin{lemma}[槽覆盖识别高支撑目标]\label{eighth:good-slots}
取 \(n\ge3\)。存在所有者坏根集 \(\mathcal B\)，使
\begin{equation}\label{eighth:bad-root}
 S(\mathcal B)\le C(D_{\mathcal N}+v^{9/2}S),
\end{equation}
而实际最大质量根所有者不在 \(\mathcal B\) 的高支撑边满足
\begin{equation}\label{eighth:good-high}
 \sum_{q(e)\ge3,\ {\rm owner}(e)\notin\mathcal B}|X_i\cap X_j|=O(v^8S).
\end{equation}
\end{lemma}
\begin{proof}
固定 \(\delta<3^{-(n-2)}/100\)。若一个胞某槽的未占部分至少为
\(\delta r_i^{n-2}\)，称其缺槽。由 \(R_i\asymp r_i\)，缺槽胞总根权至多
\(CD_{\mathcal N}\)。固定窗口 \(R_*=100(n+1)\)，令坏根为 \(Z_v\) 以及
好根中其完整 \(R_*\) 清单含缺槽胞者。
清单尺度可比和窗口 packing 重数给 \eqref{eighth:bad-root}。

先看一个不缺槽胞 \(A\) 和其真正正 \(q\) 面邻居 \(B\)。
原构造给法向中心差 \(r_A+O(vr_A)\)，每个切向中心差至多
\((1/2+1/128)r_A\)。固定切向轴 \(l\)；
两胞在其余轴的公共投影体积至少为 \(3^{-(n-2)}r_A^{n-2}\)。
故可在公共块中选一个属于 \(A\) 的 \((q+,l+)\) 正常占用的点。
该正常清洁片重标为四胞
\(\begin{smallmatrix}C&D\\ A&\widetilde B\end{smallmatrix}\)。
必须有 \(\widetilde B=B\)：若不同，二者平面交面积至少 \(r_A^2/4\)，
横向又同含此点；经 \(A\) 比较，两者都在正常边锚点的完整固定清单，
触发其自身阈值 \(O(v^3)<1/100\) 的碰撞，矛盾。

若正常对角为 \(AD\)，实际 \(l\) 端点满足
\[
 B_l^+-A_l^+=-L_{AD,l}-\Gamma_{BD}.
\]
若为 \(BC\)，则
\[
 B_l^+-A_l^+=L_{BC,l}+\Gamma_{AC}.
\]
辅助面对 \(BD\) 或 \(AC\) 在其他轴有固定交裕量；
原中心分离使其法向半盒缝非负，四胞布局使缝小于 \(r_A/8\)。
其几何切向裕量至少为 \(1/16\)，故实际调用 \(\tau_0=1/128\) 的稳健面预算。
反射 \(l\) 得左端同样结论，所有符号亦然。因此真正面一步在每个切向轴
传递左右实际端点的误差至多为 \(C T_0v^3r_A\) 加一个稳健面隙。

对好所有者 \(i\)，令 \(\mathcal E_i\) 为其 \(R_*\) 清单中全部上述稳健面隙
及 CDF 端移之和。固定窗口重复次数给
\(\sum_{i\notin\mathcal B}\mathcal E_ir_i^{n-1}=O(v^6S)\)。
取一条 \(q\ge3\) 高支撑边，所有者 \(R_i\ge R_j\)。
若 \(\mathcal E_i\ge c_*vr_i\)，原交量已不超过
\(Cv^{q-1}\mathcal E_ir_i^{n-1}\)。
否则依次沿每个临界轴的正方向走一次真正面邻居，共至多 \(n\) 步，终点为 \(h\)。
全路径及辅助胞都在该固定窗口，节点不缺槽。
稳健实际根差由 \(2\Gamma+\Delta+\Delta\) 控制，故
\[
 |R_h-R_i|\le C_n\mathcal E_i,\qquad
 |X_{h,l}^--X_{i,l}^+|\le C_n(T_0v^3r_i+\mathcal E_i)
\]
在临界轴成立；非临界轴的左右端分别接近 \(X_i\) 的对应端。

需识别实际终点 \(h=j\)。在临界轴两中心相差 \(O(vr_i)+C_n\mathcal E_i\)。
在非临界轴，原判据 \(A_1=5(1+A_8)\) 及所有者次序给
\[
 |x_{j,l}-x_{i,l}|<
 \frac{1-2A_1v}{1-A_8v}R_i,\qquad
 |x_{h,l}-x_{i,l}|
 \le(2A_8+o(1))vR_i+C_n(T_0v^3r_i+\mathcal E_i).
\]
又 \(r_h\ge R_h\ge R_i-C_n\mathcal E_i\)。
先选小 \(c_*\)，再选小 \(v\)，这些严格裕量使
\(\|x_h-x_j\|_\infty<r_h\)，若 \(h\ne j\) 则违反原强分离。
因此所有临界交长均至多 \(C_n(T_0v^3r_i+\mathcal E_i)\)。
结合原 \(L_l\le C v r_i\)，分 \(\mathcal E_i/r_i\le v^3\) 或其反面，得
\[
 I_e\le C(v^{3q}r_i^n+v^{q-1}\mathcal E_ir_i^{n-1}).
\]
有限度求和即为 \(O(v^8S)\)。
\end{proof}

对固定管宽 \(K>0\)，置
\[
 \begin{aligned}
 \mathcal C_i(K)&=\{x\in x_i+W_i':
 \text{至少两个轴到有效盒某面的距离}\le Kv^3r_i\},\quad\\
 Z(K)&=\frac va\sum_i|[(x_i+W_i')\setminus F_i]\setminus\mathcal C_i(K)|.
 \end{aligned}
\]
\begin{lemma}[一般支撑数的实际净价格]\label{eighth:high-net}
对每个固定 \(\eta,K>0\)，
\begin{equation}\label{eighth:high-payment}
 \sum_{q(e)\ge3}I_e\le
 \eta Z(K)+\eta v^6D_{\mathcal N}+C_{\eta,K}v^8S .
\end{equation}
\end{lemma}
\begin{proof}
二维没有高支撑，结论直接成立；以下取 \(n\ge3\)。
只需处理坏所有者三临界边及全部四临界以上边，其余由
引理~\ref{eighth:good-slots} 处理。
对深处 \(q\ge3\) 边，以最大实际根端为所有者，排序
\(0<p_1\le\cdots\le p_q\le C_pv\)，其中 \(p_l=L_l/r_i\)。
令 \(P=\prod_{l=2}^q p_l\)、\(A=p_1P\)、\(H=\min(p_1,A/v^q)\)；
实际交量 \(I_e\le C_IAr_i^n\)。
目标缺量例外 \(\varepsilon_j\ge c_0A/v^q\) 给
\(I_e\le Cv^q\varepsilon_jr_i^n\)。
其余若 \(\alpha_i+\Delta_i/r_i\ge c_1H\)，由
\(A/H=\max(P,v^q)\le Cv^{q-1}\) 得相应面例外界。
二者总交量为 \(O(v^{q+5}S)\)，全部合计 \(O(v^8S)\)。

非例外时在有效盒内取有向面深
\[
 [c_3H,2c_3H]r_i,\quad
 [c_2p_2,2c_2p_2]r_i,\ldots,[c_2p_q,2c_2p_q]r_i,
\]
非临界轴留 \(Kv^3r_i\) 核心。
先选小 \(c_2\)，再选小 \(c_1,c_3\)，使这些带在各有效宽四分之一内。
原始面对端距至多 \(R_i-L_l\)，端移小于 \(c_1H\)，故同一原见证
每个临界轴至少捕获 \(\min(p_lr_i/2,w^0_{j,l})\ge\gamma p_lr_i\)；
这保留真实截宽，非临界轴整宽进入。
最后取小 \(c_0\)，扣实际 \(\mathscr D_j\) 后仍有
\(car_i^nA/v^q\) 的额外 \(F_j\) 质量。
真正第1轴点态面对容量却把半径增量限制为
\(C(2c_3+c_1)Hr_i\)，剩余容量至多其固定倍 \(ar_i^{n-1}\)。
因 \(H\le A/v^q\)，再缩小 \(c_1,c_3\) 即矛盾；
所以整矩形 \(\mathcal R_e\) 缺源。

取固定 \(T(K)>2K/c_2+1\)。若 \(p_2>Tv^3\)，
第2至q轴及所有非临界轴到两面均远于管阈值，最多第1轴近面，
故矩形在整个管外。其成本
\[
 X_e=(v/a)|\mathcal R_e|
 \ge c_Xr_i^nHP/v^{q-1}
 \ge cI_ev^{1-q}\min(1,P/v^q).
\]
同一所有者跨全部支撑数至多有 \(D\) 个相交矩形，故
\(\sum_{{\rm owner}=i}X_e\le DZ_i(K)\)，只花一份物理剩余。
分配每边 \((\eta/D)X_e\)，记 \(\lambda=\eta c_X/(DC_I)\)。
在饱和分支小 \(v\) 已支付 \(A\)；非饱和而尚有正净额时
\(P<v^{2q-1}/\lambda\)。排序 \(p_1^{q-1}\le P\) 给
\[
 I_e\le(\eta/D)X_e+C_\eta v^{\beta_q}r_i^n,\qquad
 \beta_q=\frac{q(2q-1)}{q-1}=2q+1+\frac1{q-1}.
\]
于是三临界长尾净余为 \(C_\eta v^{15/2}S(\mathcal B)\)，
\(q\ge4\) 的净余晚于 \(v^8S\)。
短窗 \(p_2\le Tv^3\) 的 gross 直接至多 \(C_Kv^{q+4}r_i^n\)；
三临界只保留 \(C_Kv^7S(\mathcal B)\)，不从管外 Z 支付。
边界高支撑用原单源界和边界邻端权 \(O(v^5S)\)，为 \(O(v^8S)\)。
由 \eqref{eighth:bad-root}，坏根项均被小 \(v\) 下任意固定
\(\eta v^6D_{\mathcal N}\) 吸收，其 \(v^{9/2}S\) 余项更小。
这证明结论；仅扣原来已保留的 Z 和槽领先预算，不另减任何 \(J_i\)。
二维没有高支撑，式中左侧为零，无须越界调用该几何。
\end{proof}

\subsubsection{完整标量亏损与定量正常反馈}

令 \(f(s,t)=s_+t_+(s_++t_+)\)，并置
\[
 D(x,y,\zeta)=\tfrac12[f(x+\zeta,y+\zeta)+f(x-\zeta,y-\zeta)],
 \qquad G(b,x,y,\zeta)=b/27-bxy+D.
\]
正部凸性给 \(D\ge xy(x+y)\)，从而
\[
 G\ge b/27-\sup_{x,y\ge0}xy(b-x-y)=b(1-b^2)/27\ge0.
\]
在固定正常窗外，存在与后选 \(T,m\) 无关的 \(c_\nu>0\)，使
\begin{equation}\label{eighth:global-coercivity}
 G\ge c_\nu(b+D).
\end{equation}
证明可用反例序列：\(xy\le(D/2)^{2/3}\) 排除 \(D/b\to\infty\)，
也排除 \(b\to0\) 且 \(D/b\) 有界；余下 \(b\) 有正下界、\(D\) 有界，
近取等迫使 \(b=1,x=y=1/3,\zeta=0\)。
截幅约束使 \(\zeta\) 有界，正宽处严格凸性使最后等号唯一。
零权零成本点无需控制。

正常窗内还可取显式二次界
\begin{equation}\label{eighth:quadratic}
 G\ge10^{-3}\{(1-b)+(x-1/3)^2+(y-1/3)^2+\zeta^2\}.
\end{equation}
取 \(\nu\) 使 \(1/6\le x,y\le1\)、\(b\ge1/2\)、\(|\zeta|\le1/12\)。
正部未触零，\(\zeta\) 贡献恰为 \(3(x+y)\zeta^2\)。
余式对 \(b\) 仿射，只须检查 \(b=1/2,1\)。
前者至少为 \(1/72\)；后者置 \(s=x+y,d=x-y\)，在 \(s\le1\) 时为
\[
 \tfrac14(s-2/3)^2(s+1/3)+\tfrac14(1-s)d^2.
\]
在 \(s\le5/6\) 直接控制两平方差；在 \(5/6\le s\le1\) 第一项至少 \(7/864\)；
在 \(s\ge1\) 余式至少 \(1/27\)。均强于所写常数。

下面先给所有有界目标共用的实际强迫式。
固定 \(T\)，在同源、同有向轴对、同排序 \(s,s+y\) 上，
其他轴留共同核心，距两端大于 \(H_Tv^3R_i\)。
记
\[
 P_{ij}=(L_1+c_{ij})_+,\quad Q_{ij}=(L_2+c_{ij})_+,\quad
 c_{ij}=\operatorname{clip}((R_i-R_j)/2;-\max L_l,\max L_l),
\]
\[
 S_i^\perp=\prod_{\ell\perp}(w'_{i\ell}-2H_Tv^3R_i),\qquad
 T_j^\perp=\prod_{\ell\perp}w^0_{j\ell},\qquad
 \mathcal P_j(y)=[0,P_{ij}]\times[0,(Q_{ij}-y)_+]\times W^0_{j,\perp}.
\]
对任何被选有界目标子集，实际源点必须满足
\begin{equation}\label{eighth:shared-height}
 \begin{aligned}
 &s\ge\sum_j A_{ij}(y)-CR_iE_i(T),\qquad\\
 &A_{ij}(y)=
 \frac{T_j^\perp P_{ij}(Q_{ij}-y)_+-|\mathscr D_j\cap\mathcal P_j(y)|}
      {naR_i^{n-1}}\ge0,\\
 &\quad \sum_iE_i(T)R_i^n=O_T(v^6S).
 \end{aligned}
\end{equation}
非负性是帽内实际 \(F_j\) 质量，不能在嵌套层比较中丢掉。
若原半径增量 \(\ell\le s\)，目标帽只损失厚
\((s-\ell)+\Delta_i\) 的边带，其斜率是容量导数的 \(O_T(v)\) 倍，
故用容量在 \(s\) 的值上界即可。
若 \(\ell\ge s\)，同一个见证同时看见真正浅面邻居 \(h\) 的整横截面轴帽
\([A_h(\ell-s-g_{i\to h}-\Delta_i)-|\mathscr D_h|]_+\)。
它与自身 \(G_i\) 及组中角目标排他；法向全捕获会超过剩余容量。
原见证区间内 \(na(R_i+\ell)^{n-1}<A_h\)，容量减轴帽在转折前递增、后递减，
可用转折点 \(s+g_{i\to h}+\Delta_i+|\mathscr D_h|/A_h\) 的容量上界。
因此可取
\[
 E_i(T)=C_T\left[v^8+\alpha_i+\Delta_i/R_i+
                    \max_h|\mathscr D_h|/(A_hR_i)\right].
\]
粗先验、面邻尺度可比及有限度给所需均值界。
不同目标的实际质量先在同一个原见证中相加，
再对总高度取一个正部；不同排序和有向轴组用共同核心分离。

鲜导数乘实际层并作两排序积分，给每个有向端的完整主成本
\begin{equation}\label{eighth:nominal-cost}
 \frac{(1+v)^{n-1}}{2nv^3}\frac{R_j^{n-2}}{R_i}
                    P_{ij}Q_{ij}(P_{ij}+Q_{ij}),
\end{equation}
并扣除真实退款核
\begin{equation}\label{eighth:refund-kernel}
 \frac{(1+v)^{n-1}(v/a)S_i^\perp}{naR_i^{n-1}}
 \int_{\mathcal P_j(0)}
 (P_{ij}+Q_{ij}-x-y)\ind_{\mathscr D_j}\dd z .
\end{equation}
核由 Tonelli 的出现长度 \(Q_{ij}-y\)、\(P_{ij}-x\) 相加得到，
包含目标全部原横向切片。高度侵蚀总价格为 \(O_T(v^8S)\)，
导数 \(\varepsilon_i\) 损失为 \(O_T(v^{11}S)\)，
核心删宽为 \(O_T(v^8S)\)，宽度平均误差为 \(O_T(v^{17/2}S)\)。

仅选正常目标、用同一个核心和 \(E_i(T)\)，所得实际反馈层记为 \(Y_{i,N}\)。
它位于固定正常小管 \(\mathcal C_i(K_N)\)，\(K_N\) 不依赖 \(T\)，
每胞体积至多 \(Ca v^5R_i^n\)。
定义
\[
 R_N=\frac va\sum_i|\mathscr D_i\setminus Y_{i,N}|,\qquad
 \mathfrak E=\sum_iR_i^{n-1}\Delta_i,\qquad
 \mathfrak G_N=v^6\sum_{e\in\mathcal N}G(b,x,y,\zeta)R_e^n .
\]
置 \(p_0(v)=n/[3(1+v)]\)，标准臂为
\[
 0\le y\le e_i,\quad0\le x\le h_i(1-y/e_i),\qquad
 e_i=p_0v^3R_i,\quad h_i=p_0^2v^4R_i/n,
\]
及交换两轴，其他轴取原端整宽；各正常角之并为 \(\mathcal A_i^0\)。
若 \(d_e=|x-1/3|+|y-1/3|+|\zeta|+\sqrt{1-b}\)，则
\begin{equation}\label{eighth:quantitative-profile}
 \sum_i|\mathscr D_i\triangle\mathcal A_i^0|
 \le C(a/v)R_N+Ca v^5\sum_{\mathcal N}d_eR_e^n
                   +Ca v^6S+Cv^{n-1}\mathfrak E.
\end{equation}
先因 \(Y_{i,N}\subset\mathscr D_i\) 支付未解释缺集。
从 \eqref{eighth:shared-height} 去掉实际目标缺量，核界 \(Cv^2/a\)
与 \(\sum|\mathscr D_i|=O(av^5S)\) 给物理层误差 \(O(av^6S)\)。
剩余标称层在固定紧窗内对参数 Lipschitz，给显示的 \(d_e\) 项。
排序变量改为物理第二面距的相对误差为 \(O(v)\)；
补回核心、宽度及正常根比误差更小，端面转换另付
\(C(vR_i)^{n-1}\Delta_i\)。这证明该式，未用任何无速率形状结论。

由此，正常总退款满足
\begin{equation}\label{eighth:quantitative-refund}
 \mathfrak F_N\le
 \Gamma v^7W_N+\epsilon\mathfrak G_N+
 Cv(R_N+\mathfrak E)+C_{\epsilon,T}v^8S
\end{equation}
对任意固定 \(\epsilon>0\) 成立。
核界将 \eqref{eighth:quantitative-profile} 的误差变为
\(Cv^7\sum d_eR_e^n+Cv(R_N+\mathfrak E)+Cv^8S\)。
本角二维矩为
\[
 2\int_0^{e_j}(2e_j-z)h_j(1-z/e_j)\dd z=\tfrac53e_j^2h_j.
\]
须先乘回目标横宽 \(T_j^\perp\) 及 \eqref{eighth:refund-kernel} 前因子；
标称乘积恰为
\[
 \frac{(1+v)^{n-1}5p_0(v)^4}{3n^2}
             v^7\frac{R_j^{n+1}}{R_i}.
\]
两端主系数为 \(\Gamma\)；其他同轴反符号臂不交，共一轴或无共轴的臂
分别多 \(v^2,v^4\) 的相对因子。帽偏移和横权转换由 \(d_e\) 支付，
横宽平均误差更小。最后 \eqref{eighth:quadratic} 和有限度给
\[
 v^7\sum d_eR_e^n\le Cv^4\sqrt{S\mathfrak G_N}
 \le\epsilon\mathfrak G_N+C_\epsilon v^8S.
\]
这就是平方根误差只花八阶的原因。

正常成本与交量之差同时保留完整 \(\mathfrak G_N\)：
二者及领先预算除以 \(n^2(1+v)^{n-4}v^6R_e^n\)，
分别为 \(D,\widehat bxy,\widehat b/27\)，总误差 \(O_T(v^8S)\)。
正常根比误差是 \(O(v^3)\)；把 \(\widehat b\) 截成 \(b\)，
因窄根端的横宽上界有
\(\sum(\widehat b-1)_+R_e^n=O(v^{7/2}S)\)，成本更高阶。
故不能先只用饱和得到一个无速率 \(d_e\to0\)，再把反馈误差直接丢弃。

\subsubsection{全二轴共同层、非正常折扣及两管次序}

\begin{lemma}[全部二轴联合价格]\label{eighth:joint-price}
存在固定 \(c_R,c_D>0\)，使对任意预定 \(\chi>0\)，最后取小 \(v\) 时
\begin{equation}\label{eighth:joint}
 \begin{aligned}
 &H_2+c_RZ(K_{\rm all})+c_Dv^6D_{\mathcal N}\\
 &\le[A_6(v)+\Gamma v^7]\mathcal C+\chi\mathcal F+Cv^8S,\\
 &\quad H_2=\sum_{q(e)=2}I_e-\sum_iJ_i,\quad\\
 &A_6(v)=\frac{k}{(1+v)^2}v^6 .
 \end{aligned}
\end{equation}
其中 \(K_{\rm all}\) 是固定大管，且 \(\mathcal N,O_s,D_{\mathcal N}\)
均为上述同一正常集合及清洁分配。
\end{lemma}
\begin{proof}
先固定正常窗、其小管 \(K_N\) 及 \eqref{eighth:global-coercivity} 的常数。
再取大 \(T\)，随后取
\[
 0<m\le\min(n/12,k/(16T)),\qquad
 C_{\rm soft}\ge\max(1/m,16T/k,12/n,1).
\]
有界正常及有界非正常集合分别记为 \(\mathcal N,\mathcal B_2\)，后者
满足 \(m\le p,q\le T\)。其余有界小角为 \(\mathcal D\)，
长角 \(\mathcal L\) 满足 \(\max(p,q)>T\)，边界单列。
所有源和角保留；不按好根筛掉这些真实交量。

将 \(\mathcal N\cup\mathcal B_2\) 的全部实际目标质量代入
\eqref{eighth:shared-height}，使用同一个误差 \(E_i(T)\)、同一核心及排序，
得一个共同实际禁层 \(Y_{i,\rm all}\)。
仅留下正常目标得原 \(Y_{i,N}\)。每个 \(A_{ij}\ge0\)，因此
\[
 Y_{i,N}\subset Y_{i,\rm all}\subset(x_i+W_i')\setminus F_i.
\]
令 \(B_N,B_{\mathcal B_2}\) 为这两类目标的完整名义成本，
即在共同源核心对完整原目标横宽的高度先线性积分，再乘
\((1+v)^{n-1}v/a\)；不先压成 \(b_e\) 倍。
置 \(R_{\rm all}=(v/a)\sum|\mathscr D_i\setminus Y_{i,\rm all}|\)。
正部单调且 1-Lipschitz，精确给
\begin{equation}\label{eighth:nested-residual}
 R_N\le R_{\rm all}+(1+v)^{-(n-1)}B_{\mathcal B_2}
       \le R_{\rm all}+B_{\mathcal B_2}.
\end{equation}
此系数是一，不依赖 \(T,m\)；同角两类层不需要互不相交。
最后才选覆盖全部有界层的大管 \(K_{\rm all}\)，有
\(Y_{\rm all}\subset\mathcal C(K_{\rm all})\)、\(Z(K_{\rm all})\le R_{\rm all}\)。

只对这一个共同层使用一次 \eqref{eighth:cdf-reserve}，完整盈余给
\[
 \sum_iJ_i\ge B_N+B_{\mathcal B_2}-\mathfrak F_N-\mathfrak F_{\mathcal B_2}
                  +c_*(R_{\rm all}+\mathfrak E)-C_Tv^8S.
\]
正常退款按 \eqref{eighth:quantitative-refund} 和
\eqref{eighth:nested-residual} 控制；乘 \(v\) 的常数与 \(T\) 无关，
因为正常核窗固定。非正常退款的核为 \(C_Tv^2/a\)；
将目标缺集分成 \(Y_N\) 和其余，用 \(p,q\ge m\) 给每条完整成本
至少 \(c_mv^6R_e^n\)，得
\[
 \mathfrak F_{\mathcal B_2}
 \le C_{T,m}v(R_{\rm all}+B_{\mathcal B_2})+C_Tv^8S.
\]
非正常全域强制性同时保留两份固定余量：
\begin{equation}\label{eighth:bad-discount}
 \sum_{\mathcal B_2}I_e-B_{\mathcal B_2}
 \le(1-2\gamma_\nu)A_6W_{\mathcal B_2}
                 -c_\nu'B_{\mathcal B_2}+C_Tv^8S .
\end{equation}
这里 \(\gamma_\nu,c_\nu'>0\) 不依赖后选 \(T,m\)。
完整 \(D\) 成本与实际名义成本有固定双侧比较；
非正常根比误差是 \(Cv\) 倍完整成本，由强制性的固定一份吸收，
宽度和核心误差为绝对 \(O_T(v^{17/2}S+v^8S)\)。
将 \eqref{eighth:global-coercivity} 的 \(b+D\) 正量分成小的两份，
分别得到角权折扣和完整成本预留，不能将同一份全亏损使用两次。

小角不申请源盈余，其 gross 本身满足
\(pq\le(k/16)\omega_e\le\tfrac14 k(1+v)^{-2}\omega_e\)。
边界原单源界给
\(C[v^2\sum_{\text{邻边界}}\varepsilon_iR_i^n+v^4S(\partial I)]\)。
由 \(\mu_i\ge c aR_i^n\)、\(W_i\subset W_i^0\)，
\(\varepsilon_i aR_i^n\le C\varepsilon_i\mu_i
=C(|W_i|-\mu_i)\le C|\mathscr D_i|\)；
把缺集分成 \(Y_N\) 和余量，再用边界邻端根权 \(O(v^5S)\)，
该费至多 \(Cv(R_{\rm all}+B_{\mathcal B_2})+Cv^9S\)。

还须支付所有长而退化的二轴角，且不让 \(T\) 依赖最后的大管。
所有者取大根端，以原尺度记
\(P=\max L/r_i,Q=\min L/r_i,H=\min(Q,PQ/v^2)\)。
目标保证帽两轴宽为固定倍数 \(Pr_i,Qr_i\)，其余取原整宽。
任何正常层两轴支撑长度至多 \(C_Nv^3r_j\)，故
\begin{equation}\label{eighth:long-cap-fraction}
 |Y_{j,N}\cap\mathcal P|
 \le C(v^3/P+v^2)|\mathcal P|.
\end{equation}
若层轴对含长帽轴，由该轴取比例 \(Cv^3/P\)；
若不含，至少一层轴在帽中取完整横宽 \(\ge cvr_j\)，给 \(Cv^2\)。
此论证不除以任意小的 \(Q\)，常数不依赖 \(T\)。
先取大 \(T\) 再小 \(v\)，正常层占目标帽不超过四分之一。
若目标剩余缺量过大，全部交损由
\(CvR_N\le Cv(R_{\rm all}+B_{\mathcal B_2})\) 支付。
面隙／端移例外由 \(PQ\le CvH\) 给 \(Cv(\mathcal F+\mathfrak E)\)。

其余长角取实际缺源条，长轴深度 \([c_2P,2c_2P]r_i\)，
短轴 \([c_3H,2c_3H]r_i\)，其他轴留正常小管 \(K_N\) 的核心。
同一原见证真实截宽后的目标量仍至少 \(car_i^nPQ/v^2\)，
浅轴容量至多固定小倍数 \(ar_i^nH\)，故整条缺源。
其成本至少
\[
 (v/a)|\mathcal R_e|\ge
 cr_i^nPQ\min(1/v,P/v^3)\ge cT I_e .
\]
取 \(T>4K_N/c_2\)，该条避开正常小管及 \(Y_N\)。
按源度数分配，长角总费至多
\begin{equation}\label{eighth:two-tubes}
 (C_L/T)R_N\le(C_L/T)(R_{\rm all}+B_{\mathcal B_2}).
\end{equation}
穿透比例、正常支撑常数、核心和度数均独立于 \(T,m,K_{\rm all}\)，
所以 \(C_L\) 与后选 \(T\) 无关。
长条无需避开覆盖全部有界层的大管：与非正常层重合的部分已经由
\eqref{eighth:nested-residual} 的完整成本预留支付。

现在先以小固定 \(\epsilon\) 吸收正常退款的 \(\mathfrak G_N\)；
选固定 \(T\) 使 \(C_L/T\) 小于 \(c_*,c_\nu'\) 的预留份额，
最后令 \(v\) 小吸收全部 \(C_{T,m}v\) 项。
剩下固定的 \(c_RR_{\rm all}\)，并得
\[
 H_2+c_RZ(K_{\rm all})
 \le[A_6+\Gamma v^7]W_N+(1-\gamma_\nu)A_6W_O+Cv\mathcal F+Cv^8S .
\]
结合 \eqref{eighth:slot-identities}，价格右端至多
\[
 [A_6+\Gamma v^7]\mathcal C-
       [\gamma_\nu A_6+\Gamma v^7]D_{\mathcal N}+O(v^9S).
\]
这给固定的 \(c_Dv^6D_{\mathcal N}\)，只用已预留的角权折扣，
不重复消费付非正常层的完整成本份额。
最后使 \(Cv\le\chi\) 即得结论。
二维空积、唯一角位及无高支撑的情形逐式成立；
正常帽和长帽估计只余同一轴对，不需越界引用高维命题。
全部局部参数按正常窗、小管、全域强制性、大 \(T\)、\(m,C_{\rm soft}\)、
共同层与最后大管的顺序固定，最后才取小 \(v\)，没有两管阈值循环。
\end{proof}

\subsubsection{一次恢复与八层粗余项}

\begin{proposition}[八阶粗上包]\label{eighth:bootstrap-upper}
对每个固定 \(n\ge2\)，有 \(C_n(1-t)\le P_{7,n}(t)+O_n(av^8)\)。
\end{proposition}
\begin{proof}
先在当前同一八层族的每个 \(m=7,8\) 末层分别应用上述量化论证。
两层各自使用本层的正常集合、槽空缺及共同层，不将第八层对象套到第七层。
只需引理~\ref{eighth:priors} 的粗先验；
特别无需先证明 \(H_m/(v^6S)\to K_n\)、\(\mathcal F_m=o(v^6S)\)
或一个无速率缺源轮廓。有限两层的局部阈值可以在统一 \(A_8\) 和 CDF 后同时固定。

在引理~\ref{eighth:joint-price} 中先固定
\(\eta<\min(c_R,c_D)/4\)，用 \eqref{eighth:high-payment}
支付全部高支撑，留正剩余得
\[
 H_m\le[A_6+\Gamma v^7]\mathcal C_m+\chi\mathcal F_m+Cv^8S .
\]
这里单临界无正体积交，故完整 \(H_m=H_2+\sum_{q(e)\ge3}I_e\)，
全部外部 pair 及完整盈余各只计一次。
二维高支撑为空，正常 packet 的横积、核心均为空积；
四槽容量由实际二维四胞论证取得，全部二轴共同价格直接逐式成立。
因此没有从只为高维表述的几何声明越界推论二维。

完整实际槽容量的乘积和由宽度先验给
\[
 \mathcal C_m=\sum_i\sum_{q<r}R_i^2
                  \prod_{\ell\ne q,r}(R_i+w_{i\ell})
 =B_n(1+v)^{n-2}\sum_iR_i^n+O(v^{7/2}S).
\]
精确质量式和 \(\Xi=O(v^6S)\) 给
\(\sum_iR_i^n=(M'+\mathcal D_m)/a\le\mathcal V/A_v+O(v^6S)\)。
于是
\begin{equation}\label{eighth:full-local}
 H_m\le Q(v)\mathcal V+\chi\mathcal F_m+Cv^8S,\qquad
 Q(v)=K_nv^6[(1+v)^{-4}+(10/9)v(1+v)^{-2}].
\end{equation}
槽容量误差乘主价为 \(O(v^9S)\)。
此处使用完整槽的实际乘积，不把软接触权上界反向当成槽容量上界。

每层只恢复一次原免费面条，正体积至少 \(c_{\rm face}\mathcal F_m\)。
取 \(\chi<c_{\rm face}/8\)。
实际小源回收的盒内 dyadic 系数保留为 \(a\)，对该层有
\[
 (v/a)\sum_{\rm deep}\delta_{i,\rm in}^{\rm mass}
 \le Ca v\sum_i\varepsilon_i r_i^n+Cv^n\mathcal F_m
 \le Cv^{n+6}S+Cv^n\mathcal F_m .
\]
盒外用同一个入向严格见证及半空间质量连续性，成本至多
\(Cv\mathcal F_m\)；边界局部小源界给
\[
 (v/a)\sum_{\partial I}\delta_i^{\rm mass}
 \le CavS(\partial I)=O(v^{n+6}S)\text{。}
\]
环境另有 \(Ca(\mathcal F_m+P_\partial)\)，其 \(aP_\partial\) 项也是
\(O(v^{n+6}S)\)。固定突出边界价格仍留给图回收。
最后取小 \(v\)，使 \(\chi\mathcal F_m\)、\(Cv\mathcal F_m\)、
\(Cv^n\mathcal F_m\)、\(Ca\mathcal F_m\) 各至多占同次收入的八分之一。
合计只用半份恢复，不重复记入另一份免费面。
\(\mathcal F_m=0\) 时这些成本也全为零。
所有 \(n\ge2\) 的其余 \(v^{n+6}S\) 成本均为 \(O(v^8S)\)。

在恢复后的实际并集上，Bonferroni 使用完整 \(H_m\)，再用
\[
 (R_i+\mu_i^{1/n})^n\ge A_vg_i/a-C(v/a)(g_i-\mu_i)
\]
及精确质量恒等式，扣掉物理 \(\mathcal D_m\) 和环境量。
重复量净系数至少 \(1/(2a)\)，小源和环境已由同次恢复支付。
因此这两个末层都可选择并集误差上包
\begin{equation}\label{eighth:two-prefix-prices}
 E_m=Q(v)\mathcal V+e_m,\qquad 0\le e_m\le C_mv^8S,\qquad m=7,8.
\end{equation}
未删除源缺量 \(\mathfrak d_m^0\) 仍保留到望远镜。
第六前缀只使用引理~\ref{sharp-source-price} 的领先价，取
\(E_6=[K_n+o(1)]v^6S\)；前五层取 \(E_h=C_0v^6S\)。
因此第七层只须重放上述量化价格，不必另加定性饱和假设。

对每层加入同一标记半盒和内部突出集，
\eqref{eq:global-recovered-union}--\eqref{eq:global-exact-interior-identity}
给同一公共左侧
\begin{equation}\label{eighth:ledger}
 \mathcal V\ge A_vM/a+c_{\rm rec}P_I-C_{\rm b}P_\partial-2nvS_I
 +c_D'\mathcal D_h/a-b(v/a)\mathfrak d_h^0-E_h,\qquad c_D'=1/2.
\end{equation}
输入只计一次。内部精确式
\[
 (1-a)S_I+M_I-A_vM_I/a
 =(1-a)(S_I-M_I/a)-(A_v-1)M_I/a
\]
与共享质量、突出量给正回收及所写体积价格；
源常数产生的 \(CvP_\partial\) 在最后小 \(v\) 时由纯维数预留吸收。

置 \(\tau=bv/c_D'\)、\(\Omega_8=\sum_{j=0}^7\tau^j\)。
使用 \eqref{eighth:depth} 以 \(\tau^{8-h}\) 加权，精确有
\begin{equation}\label{eighth:telescope}
 \sum_{h=1}^8\tau^{8-h}(c_D'f_h-bvf_{h-1})
                  =c_D'f_8-bv\tau^7f_0 .
\end{equation}
初端除 \(a\Omega_8\) 为
\(A_0b^8v^9(M'/a)/((c_D')^7\Omega_8)=O(v^9S)\)。
对 \(\bar E=\Omega_8^{-1}\sum_h\tau^{8-h}E_h\)，先减同一 \(Q(v)\mathcal V\)：
\[
 \bar E-Q(v)\mathcal V=
 \frac{e_8+\tau e_7+\tau^2[E_6-Q(v)\mathcal V]
       +\sum_{h=1}^5\tau^{8-h}[E_h-Q(v)\mathcal V]}{\Omega_8}.
\]
四类项分别为 \(O(v^8S)\)、\(O(v^9S)\)、\(o(v^8S)\)、\(O(v^9S)\)；
第三项只需 \eqref{eighth:leading-mass} 的领先比较。
图条件吸收边界及内部体积价格，\(f_8\) 和剩余回收非负，遂得
\[
 \mathcal V\ge A_vM/a-Q(v)\mathcal V-Cv^8S-o(v^8S).
\]
零外部按零源口径，不定义空 CDF，内部计数给同方向。

由 \(S/V\le(1-a)^{-1}\)，直接移真实 \(M\) 得
\[
 1+M/V\le\frac{A_v}{A_v-a(1+Q(v))}+O(av^8)+o(av^8).
\]
与 \(L_n=A_v/(A_v-a)\) 的精确差为
\[
 \frac{aA_vQ}{(A_v-a)(A_v-a-aQ)}.
\]
因 \(a=O(v^2)\)，只须
\[
 Q=K_nv^6-\tfrac{26}{9}K_nv^7+O(v^8),\qquad
 A_v/(A_v-a)^2=1-nv+O_n(v^2),
\]
即得 \(P_{7,n}+O_n(av^8)\)。
若全部输入 big-O 失败，可按上确界取输入使归一化余量趋于 \(+\infty\)，
再以 \(o(Vav^8)\) 的精度取外部紧集；该序列最终进入
\eqref{eighth:critical}，与上述固定常数加逐序列小量矛盾。
所有原正半径及尺度跨度均保留，不将固定输入的 little-o 与上确界交换。
\end{proof}

\subsubsection{竞争输入的正预算稳定性}

以下进一步限制于任意固定 \(B<\infty\) 的序列
\begin{equation}\label{eighth:exact-competitive}
 1+M/V\ge P_{7,n}(t)-Bav^8,\qquad v\downarrow0.
\end{equation}
同一 \(I\)、八层、原指定见证和统一 CDF 均保持不变。
正常窗、清洁槽分配和前面共同层的固定阈值也先固定。
随后出现的辅助有限窗只用于测试新的目标 packet。

\begin{lemma}[八阶稳定预算]\label{eighth:exact-stability}
在第八层，可固定常数 \(C_*\)，使
\begin{equation}\label{eighth:exact-stable-budget}
 \mathfrak G_N+B_{\mathcal B_2}+R_{\rm all}+R_N+\mathfrak E
 +\mathcal F+\mathcal F_{1/128}+v^6D_{\mathcal N}
 +P_I+P_\partial\le C_*v^8S.
\end{equation}
其中正常层和 \(\mathcal F_{1/128}\) 均取本层的原实际对象；
\(\mathcal F_{1/128}\) 包含固定窗口中全部具有该切向裕量的稳健面对。
具体沿用 \(\mathscr E_\tau\) 的稳健近面定义，
\(\mathcal F_\tau=\sum_{(i,p)\in\mathscr E_\tau}\Gamma_{ip}r_i^{n-1}\)；
固定窗口及有界重复次数均在八层统一常数之前取定。
特别
\[
 \sum_{\mathcal N}d_e^2R_e^n=O(v^2S),\qquad
 D_{\mathcal N}=O(v^2S).
\]
常数不依以后选取的大角窗或尾份额。
\end{lemma}
\begin{proof}
重开引理~\ref{eighth:joint-price} 的同一次证明。
\eqref{eighth:quantitative-refund} 的平方根误差用 Young 只花
\(\mathfrak G_N\) 的一个固定小份额。\eqref{eighth:global-coercivity}
的 \(b+D\) 另分固定份额，分别保留完整 \(B_{\mathcal B_2}\)、
非正常角权折扣和退款吸收。长条的 \(C_L/T\) 先取小，
其余 \(C_{T,m}v\) 最后取小；由
\eqref{eighth:nested-residual}，这仍保留固定正系数的
\(R_{\rm all}+\mathfrak E+v^6D_{\mathcal N}\)。
故对固定正 \(c_j\)，原联合证明给
\[
 H_2+c_1\mathfrak G_N+c_2B_{\mathcal B_2}
 +c_3R_{\rm all}+c_4\mathfrak E+c_5v^6D_{\mathcal N}
 \le[A_6+\Gamma v^7]\mathcal C+Cv\mathcal F+Cv^8S.
\]
先固定 \(\eta_0<\min(c_3,c_5)/16\)，再用
引理~\ref{eighth:high-net}，高支撑只占其固定小部分。
因 \(Z(K_{\rm all})\le R_{\rm all}\)，全部正量仍有固定份额。

普通面条和固定稳健面条各有实际正体积下界。
二者即使重叠，其并集至少为两个下界的各一半之和；
有限多个候选组也可用固定正权平均。条均避开 \(\bigcup X_i\)
及标记部分，位于原源的 \(Q(x,R_i)\) 内，环境仍由原固定面邻居控制。
于是同一次恢复可保留正的 \(\mathcal F+\mathcal F_{1/128}\)。
全部 \(Cv\mathcal F,Cv^n\mathcal F,Ca\mathcal F\) 支出合并，
最后小 \(v\) 时只占其固定小份，不能逐项重用整份面收入。

将这些正量保留在第八层的 \eqref{eighth:ledger}，以原八层权平均，
末层权 \(1/\Omega_8=1+O(v)\)，不损失固定正系数。
每层公共的内部收入按 \eqref{sharp-graph-absorption} 支付边界和
\(2nvS_I\) 后仍留固定正的 \(P_I\)；内部精确式的 \(1-a\) 不变。
故以 \(\mathcal L\) 记所列源／面正量，有
\[
 (1+Q)\mathcal V-A_vM/a
 \ge c_PP_I+c_L\mathcal L-Cv^8S-o(v^8S).
\]
\eqref{eighth:exact-competitive}、真实比值恒等式及
\eqref{eighth:leading-mass} 给公共左侧 \(O_B(v^8S)\)。
所以 \(\mathcal L+P_I=O(v^8S)\)，原图再给 \(P_\partial=O(v^8S)\)。
最后用 \eqref{eighth:nested-residual} 得 \(R_N\) 的界，
用 \eqref{eighth:quadratic} 得参数平方界。
本证明只使用已证粗上包，没有使用下面的新八阶价格。
\end{proof}

\begin{lemma}[全部高支撑交量的小量]\label{eighth:exact-high-gross}
对上述同一族，
\[
 \sum_{q(e)\ge3}I_e=o(v^8S).
\]
二维为空和；结论不再扣除任何源盈余。
\end{lemma}
\begin{proof}
先用实际质量和旧正常层的逐源体积上界得到
\[
 \varepsilon_iR_i^n\le Cv^5R_i^n+
       C|\mathscr D_i\setminus Y_{i,N}|/a.
\]
对任一实际指标集 \(A\)，有限相交度和根比给
\[
 \sum_{\operatorname{Adj}(A)}\varepsilon_iR_i^n
 \le Cv^5S(A)+(C/v)R_N.
\]
引理~\ref{eighth:good-slots} 的真实缺槽坏根集 \(\mathcal B\)
满足 \(S(\mathcal B)=O(v^2S)\)；边界权为 \(O(v^5S)\)。
两者的邻端 epsilon 权均为 \(O(v^7S)\)。

好所有者的固定面对路径有切向裕量 \(1/16\)，故属于
预先取定的 \(1/128\) 稳健候选。其路径长度误差均值现在为
\(O(\mathcal F_{1/128}+\mathfrak E)=O(v^8S)\)。
\eqref{eighth:good-high} 因而给总交量 \(O(v^9S)\)。
坏所有者重开引理~\ref{eighth:high-net} 的实际矩形证明：
目标质量例外、面对／端移例外现在分别为
\(O(v^{q+7}S)\)，最早十阶。
其余矩形在 \(p_2>T_Kv^3\) 时避开整个原固定两轴管，
实际成本仍为 \(cr_i^nHP/v^{q-1}\)，所有支撑数共同重数有界。
所以任给固定 \(\eta>0\)，长窗总量至多
\[
 \eta Z(K_{\rm all})+
 C_\eta\sum_{q=3}^n v^{q(2q-1)/(q-1)}S(\mathcal B)
 \le\eta Z(K_{\rm all})+O_\eta(v^{19/2}S).
\]
短窗 \(p_2\le T_Kv^3\) 的逐根 gross 为
\(O_K(v^{q+4}r_i^n)\)，只乘坏根权后为 \(O_K(v^9S)\)。
边界原单源界给
\[
 O(v^{q+7}S)+O(v^{q+q/(q-1)+5}S)=o(v^8S).
\]
这里真实截宽、完整自身目标和同一原半径均由原矩形证明保留，
没有以旧 \(O(v^8S)\) 总式改名小量。
固定稳定构造先给 \(Z(K_{\rm all})\le C_*v^8S\)，
其常数只依此前固定的 \(\eta_0\)，与新 \(\eta\) 无关。
先 \(v\to0\)，再 \(\eta\downarrow0\)，即得结论。
\end{proof}

\subsubsection{薄轴帽、整源坐标与两次实际反馈}

以后固定任意有限 \(T\)，选取两端深处、两交长不超过
\(Tv^3R_e\) 的目标 packet \(\mathscr P\supset\mathcal N\)。
所有帽、根、盒和目标清单固定；测试缺集的变化不重定义几何。
记 \(\eta_h=|\mathscr D_h\setminus Y_{h,N}|\)。

\begin{lemma}[只扣 raw 的轴帽]\label{eighth:exact-axial}
在深处源 \(i\) 的有效浅面深度 \(s\) 上，同一个原指定见证使
所有 packet 帽的实际质量先相加后满足
\begin{equation}\label{eighth:exact-height-nec}
 s\ge\frac{\sum_{j\in\mathscr P}|F_j\cap\Pi_{ij}(y)|}
                  {naR_i^{n-1}}-C_Te_i,
 \qquad
 e_i=v^8R_i+\Delta_i+
          \max_{h\ \text{面邻}}(g_{i\to h}+\eta_h/\beta_h),
\end{equation}
其中 \(\beta_h\asymp(vR_h)^{n-1}\)，且
\[
 \sum_iR_i^{n-1}e_i
 \le C[v^8S+R_N+\mathfrak E+\mathcal F]=O(v^8S).
\]
该式也适用于全局最浅轴不在某个目标轴对的点。
\end{lemma}
\begin{proof}
真正面邻居 \(h\) 的原始整横截面帽 \(C_h(d)\) 中，
旧正常层占比至多
\[
 \theta_N(v)=\frac{2(n-1)K_N}{c_w}v^2+
        \frac{4\binom{n-1}{2}K_N^2}{c_w^2}v^4.
\]
轴对含帽轴时，另一短轴在原整宽中只占 \(O(v^2)\)；
不含时两条短轴各给此因子。有效面位移只移动短区间，
不增加其长度。故此比例对任意小 \(d\) 一致，令
\(\beta_h=(1-\theta_N)\prod_{\ell\ne q}w^0_{h\ell}\)，有
\[
 |F_h\cap C_h(d)|\ge
 [\beta_h\min(d,w^0_{hq})-\eta_h]_+.
\]
完整横截面的捕获由真正面邻居的固定切向裕量保证，
对所有源点及原半径均成立。若保证深度达到全宽，整个 \(F_h\)
与 \(G_i\) 同时进入见证，其质量为 \(caR_i^n\)，超过
\(O(avR_i^n)\) 剩余容量；所以实际见证不在饱和分支。

写 \(\ell=r(x)-R_i\)、\(b=s+g_{i\to h}+\Delta_i\)。
原始／有效面对端精确距离为
\(R_i+s+g_{i\to h}+\delta_i^+-\delta_h^-\)；
目标端移只改善捕获。轴帽保证深度至少 \((\ell-b)_+\)。
容量 \(C_i(\ell)=a[(R_i+\ell)^n-R_i^n]\) 在整个原半径区间的
导数严格小于 \(\beta_h\)。因此容量减轴帽在转折
\(z=b+\eta_h/\beta_h\) 前增后减，最多为 \(C_i(z)\)；
\(z\) 超过原半径区间时仍由容量单调性成立。
正常缺层只降低轴帽斜率，平移只扣 \(\eta_h\)。
原宽和 \(\eta_h\le|B_h^0|\) 还给 \(z=O_T(vR_i)\)，
故比较 \(C_i(z)\) 与 \(C_i(s)\) 时容量导数仍为 \(O(aR_i^{n-1})\)。

角帽从中点锚移到源有效端至多损 \(\Delta_i\)，移到目标原端只改善。
若 \(\ell\le s\)，总帽边带损失至多
\(D_{i,T}(s-\ell+\Delta_i)\)，其中
\(D_{i,T}\le C_Tv^{n+1}R_i^{n-1}<naR_i^{n-1}\)；
容量斜率比较给总实际帽质量至多 \(C_i(s)+D_{i,T}\Delta_i\)。
若 \(\ell\ge s\)，先将全部角目标、真正浅面邻居和自身 \(G_i\)
在同一见证中相加，再用上述转折点，给
\[
 \sum_j|F_j\cap\Pi_{ij}(y)|
 \le C_i(s)+CaR_i^{n-1}(g_{i\to h}+\Delta_i+\eta_h/\beta_h)
                    +D_{i,T}\Delta_i.
\]
各 \(F_j\) 排他，面邻居不同于二轴目标。
高度窗 \(s=O_T(v^4R_i)\) 中
\(C_i(s)/(naR_i^{n-1})=s+O_T(v^8R_i)\)，得到必要式。
有限面逆度给
\(\sum R_i^{n-1}\max_h\eta_h/\beta_h\le Cv^{1-n}\sum\eta_h=CR_N\)，
其余项由稳定预算控制。
\end{proof}

在有效盒 \(B_i'=x_i+W_i'\) 中，先选各轴最近面符号 \(\sigma\)，
再选全局最浅轴 \(q\)，平局按固定顺序。令
\[
 x_h=f_{ih}^{\sigma_h}-\sigma_h(s+y_h),\quad y_q=0,\qquad
 \Omega_{i,q,\sigma}=\prod_{h\ne q}[0,w'_{ih}/2],\quad
 H_i(y)=\min_h(w'_{ih}/2-y_h).
\]
这些 Jacobian 为一的 chart 覆盖整个有效盒至零测分界，
并保留所有三近面。仅取临界朝向与 \(\sigma\) 匹配的目标；
反向面的深度至少半个源宽，不产生本高度窗的正标称帽。
\(\Pi_{ij}(y)\) 在原始目标盒的两临界宽为
\[
 \min(w^0_{jl},(P_{ij,l}-y_l)_+),\qquad
 \min(w^0_{jm},(P_{ij,m}-y_m)_+),
\]
其余取原始整宽，其中 \(P_{ij,h}=(L_h+c_{ij})_+\)。
真实全宽截断在定义中保留；固定 \(T\) 后小 \(v\) 时自动不活跃。
令
\[
 h_i^0(y)=\frac{\sum_j|\Pi_{ij}(y)|}{naR_i^{n-1}},\qquad
 \mathcal K_i(E;y)=
       \frac{\sum_j|E_j\cap\Pi_{ij}(y)|}{naR_i^{n-1}},
\]
\begin{equation}\label{eighth:exact-operator}
 \mathscr B_{\mathscr P}(E)_i=
 \bigcup_{q,\sigma}
 \{x(s,y):0\le s<\min(H_i(y),[h_i^0(y)-\mathcal K_i(E;y)]_+)\}.
\end{equation}
输入 \(E_j\) 始终在原始目标盒，输出是其中的有效源子集。
Tonelli 保证可测性；测试 \(E\) 不改变根、锚、clip 或目标清单。

\begin{lemma}[整源小算子与两次反馈]\label{eighth:exact-iteration}
对任意测试集合族，
\[
 E\subset F\Longrightarrow\mathscr B_{\mathscr P}(F)
       \subset\mathscr B_{\mathscr P}(E),\qquad
 \|\mathscr B_{\mathscr P}(E)-\mathscr B_{\mathscr P}(F)\|_1
       \le C_Tv\|E-F\|_1.
\]
取 \(Y_0=\mathscr B_{\mathscr P}(0)\)、\(Y_1=\mathscr B_{\mathscr P}(Y_0)\)、
\(Y_2=\mathscr B_{\mathscr P}(Y_1)\)。存在真实保证层
\(B^-\subset\mathscr B_{\mathscr P}(\mathscr D)\cap D'\)，使
\begin{equation}\label{eighth:exact-second-cost}
 \sum_iJ_i\ge(1+v)^{n-1}\frac va\sum_i|Y_{2,i}|
 +c\frac va\sum_i|\mathscr D_i\setminus B_i^-|
 +c\mathfrak E-O_T(v^9S).
\end{equation}
退款的实际缺集可换成 \(Y_1\)，价格误差 \(O_T(v^9S)\)。
\end{lemma}
\begin{proof}
正部及 \(u\mapsto\min(H,u)\) 为 1-Lipschitz。
打开 Tonelli 固定一个目标物理点：若 \(q\) 在其轴对内，
可覆盖它的源 base 测度至多 \(C_T(v^3R_i)(vR_i)^{n-2}\)；
否则为 \(C_T(v^3R_i)^2(vR_i)^{n-3}\)。
各除以本源的 \(naR_i^{n-1}\)，分别为 \(C_Tv,C_Tv^3\)。
有限逆度给全族物理 \(L^1\) 小常数，二维后一类为空。
名义高度积分给 \(|Y_0|_1\le C_Tav^5S\)。

旧正常实际层 \(Y_{i,N}\) 的浅轴深为 \(O(v^4R_i)\)，伴随轴为
\(O(v^3R_i)\)，其他轴在旧核心内。它在新 chart 中仍是同一排序、
同一原始目标帽；新高度包含旧全部非负实际质量且未扣旧侵蚀。
因此 \(Y_{i,N}\subset\mathscr B_{\mathscr P}(\mathscr D)_i\) 精确成立。
\eqref{eighth:exact-height-nec} 的侵蚀 \(C_Te_i\) 只在正帽 base 有损失，
该 base 测度至多 \(C_Tv^{n+1}R_i^{n-1}\)，故物理总损失
\(e=O_T(av^9S)\)。减此侵蚀得 \(B^-\subset D'\)，并有
\[
 \rho:=\|\mathscr D-\mathscr B_{\mathscr P}(\mathscr D)\|_1
 \le (a/v)R_N+e=O_T(av^7S).
\]
未使用源取空算子及空旧正常层，原 raw 仍在该全族界内。

记 \(\lambda=C_Tv\)、\(M_0=|\mathscr D|_1=O(av^5S)\)。
反单调性给 \(Y_1\subset Y_2\subset Y_0\)，且
\[
 \|\mathscr D-Y_0\|_1\le\rho+\lambda M_0=O_T(av^6S),\quad
 \|\mathscr D-Y_1\|_1\le\rho+\lambda(\rho+\lambda M_0)=O_T(av^7S).
\]
故 \(\|\mathscr B_{\mathscr P}(\mathscr D)-Y_2\|_1=O_T(av^8S)\)。
乘源导数 \(Cv/a\) 给成本误差 \(O_T(v^9S)\)；
退款核为 \(C_Tv^2/a\)，故换成 \(Y_1\) 也只损此阶。
这里不能将全部 \(\mathscr D\) 换成 \(Y_2\) 后丢掉 raw 余量。
对 \(B^-\) 使用 \eqref{eighth:cdf-reserve} 一次：
侵蚀成本为 \(O_T(v^{10}S)\)，导数 epsilon 损失至多
\(C_Tv^6\sum\varepsilon_iR_i^n=O_T(v^{11}S)\)。
层成本、raw 和端移由同一份完整 \(J\) 同时分配，得到所述式。
\end{proof}

\subsubsection{正常 packet 的实际八阶积分}

本小节取 \(\mathscr P=\mathcal N\)，给相应迭代加下标 \(N\)。
记 \(s_v=(1+v)^{n-1}\)、\(\lambda_v=s_vv/a\)，并定义未扣任何
实际盈余的模型价格
\[
 \mathsf P_N=\sum_{e\in\mathcal N}I_e-\lambda_v\sum_i|Y_{2,N,i}|.
\]
这与只取正常交量后再减全体 \(J_i\) 不同。
令
\[
 Q_4(p,q)=\frac{pq(p^2+3pq+q^2)}3,\qquad
 Q_5(p,q)=\frac{pq(p+q)(5p^2+29pq+5q^2)}{24},
\]
\[
 F(p,q,z)=\tfrac12\{(p+z)(q+z)(p+q+2z)
                         +(p-z)(q-z)(p+q-2z)\}.
\]

\begin{lemma}[实际源成本及正常模型价]\label{eighth:exact-normal-price}
沿 \eqref{eighth:exact-competitive}，
\begin{equation}\label{eighth:exact-normal-integral}
 \begin{aligned}
 \lambda_v\sum_i|Y_{2,N,i}|
 =\sum_{e\in\mathcal N}s_vR_e^n\biggl[
 &\frac{v^6}{n}F(p,q,z)-\frac{2v^7}{n^2}Q_4(p,q)\\
 &+\frac{2v^8}{n^3}Q_5(p,q)
       +\frac{n-2}{n}v^8p^2q^2\biggr]+o(v^8S).
 \end{aligned}
\end{equation}
并有
\begin{equation}\label{eighth:exact-model-price}
 \mathsf P_N\le q_N(v)W_N+o(v^8S),\qquad
 q_N(v)=\frac{k v^6}{(1+v)^2}
       +\frac{\Gamma v^7}{(1+v)^3}+k c_{8,n}v^8,
\end{equation}
\[
 c_{8,n}=-\frac{13}{18}-\frac{n(n-2)}3+\frac{400}{243}.
\]
所有三近面源项都已包含；二维相应项为空。
\end{lemma}
\begin{proof}
先保留真实源、目标整横宽。对全局最浅轴在目标轴对内的两个排序，
名义成本为
\[
 \frac{\lambda_v S_i^\perp T_j^\perp}{2naR_i^{n-1}}
          P_{ij}Q_{ij}(P_{ij}+Q_{ij}),\qquad
 S_i^\perp=\prod_{\ell\perp}w'_{i\ell},\quad
 T_j^\perp=\prod_{\ell\perp}w^0_{j\ell}.
\]
最浅轴 \(q_0\) 不在该轴对内时，两个临界底变量独立，单向名义成本为
\[
 \frac{\lambda_vT_j^\perp}{naR_i^{n-1}}
 \sum_{q_0\perp}2\prod_{h\notin\{l,m,q_0\}}w'_{ih}\,
                         \frac{P_{ij}^2Q_{ij}^2}{4}.
\]
系数二来自该最浅轴的两个最近面。标称宽和正常根比下，
两有向端的后式正是
\[
 (n-2)s_vv^8R_e^np^2q^2/n\text{。}
\]
它的退款核为 \(O(v^3)\)，乘
\(|Y_1|_1=O(av^5S)\) 及源导数后仅为 \(O(v^9S)\)，
故名义三近面成本必须保留，三近面退款可以入余项。

整个源纤维上截断不能直接删去。某名义目标非零时，其临界底变量
为 \(O(v^3R_i)\)；总高度为 \(O(v^4R_i)\)，所以顶端截断只可能来自
非临界自由底变量靠近源中分面的宽 \(O(v^4R_i)\) 边带。
该名义项独立于该变量，相对积分损失为 \(O(v^3)\)。
乘六阶名义成本后为 \(O(v^9S)\)，多个目标相加也有同一高度上界。
无截断时高度对目标质量线性，且任意测试缺集都是原盒子集，
所以近零帽也没有负高度或相对展开问题。

固定一条正常帽。在目标的中央横截面上，同角正常目标唯一；
其他轴对需要额外坐标近面，同轴反符号角则有完整源宽间隔。
非中央横向区域中，距某面 \(O(v^4R_j)\) 的比例为 \(O(v^3)\)；
其余近面带的比例为 \(O(v^2)\)，而 \(Y_1\subset Y_0\) 再强迫
帽内某法向深度为 \(O(v^4R_j)\)，相对帽宽多 \(O(v)\)。
故非中央 \(Y_1\) 帽质量为 \(O(av^7R_j^n)\)，退款仅为 \(O(v^9R_j^n)\)。
中央 \(Y_1\) 所见 \(Y_0\) 的其他角也只有上述质量；
经一次高度分母及中央底面积作用，改变 \(Y_1\) 帽质量为
\(O(av^8R_j^n)\)，再次反馈为十阶。因此八阶退款可用同一边的
两个真实不同帽计算，不能将两个方向先设为相同。

这里的宽度替换按 \(L^1\) 而非逐源小误差进行。
平面分母 \(naR_i^{n-1}/T_j^\perp\) 除以 \(v^2R_e\)
在固定正紧区间；子图的 1-Lipschitz 性给宽度变化的物理误差
\(Cav^5(\xi_i+\xi_j)R_e^n\)。
名义成本和退款的总宽度误差分别为
\(Cv^6\sum\xi_iR_i^n=O(v^{17/2}S)\) 和更小量。
正常根比的系数替换为 \(O(v^9S)\)，但帽中的实际根差仍保留。
目标有效锚移回原始锚另付端移：每张图高度随底变量单调不增，
子图周长至多 \(C(vR_i)^{n-1}\)，有限线性 chart 变换只改常数。
故平移物理误差至多 \(C\sum(vR_i)^{n-1}\Delta_i\)，
退款为 \(Cv\mathfrak E=O(v^9S)\)。

现给双帽所需的一致截面计算。取尺度 \(e\)，\(P/e,Q/e\)
在固定正紧区间，\(|c|\le\nu e\)、\(e/f\ll1\)，
令 \(P_\pm=P\pm c,Q_\pm=Q\pm c\)。
方向 \(t\) 的首轮平面层 \(U_0^t\) 在排序 \(X\le Y\) 中满足
\(X\le P_t(Q_t-D)_+/f,\ D=Y-X\)，另一排序交换两轴。
下一层 \(U_1^t\) 的高度扣除相反方向 \(U_0^s\)
在 \([0,P_t]\times[0,(Q_t-D)_+]\) 中的实际帽质量再除 \(f\)。
这些集合保留真实排序和正部。记
\[
 M_{1,\pm}(c)=
 \int_{U_1^\mp\cap([0,P_\pm]\times[0,Q_\pm])}
           (P_\pm+Q_\pm-X-Y)\dd X\dd Y .
\]
则
\begin{equation}\label{eighth:exact-cross-cap}
 M_{1,+}(c)+M_{1,-}(c)
 =\frac{2Q_4(P,Q)}f-\frac{2Q_5(P,Q)}{f^2}
 +O\!\left(\frac{e^2c^2}{f}+\frac{e^4|c|}{f^2}
                         +\frac{e^6}{f^3}\right).
\end{equation}
为核一致性，首轮面积和自矩精确为
\[
 |U_0(P,Q)|=\frac{PQ(P+Q)}{2f},\qquad
 M_0=\frac{Q_4(P,Q)}f-\frac{P^2Q^2(P+Q)}{3f^2}.
\]
暂去外帽裁剪，双向和在 \(c\mapsto-c\) 下不变，是二阶导数
有界于 \(Ce^2/f\) 的光滑多项式。较小外帽只切长臂远端；
距端距离 \(d\) 的截厚至多 \(C(e/f)d\)，带核 \(O(e)\)
积分 \(0\le d\le2|c|\) 给 \(Ce^2c^2/f\)。
这包括 \(c\) 远大于、可比于或小于薄高 \(e^2/f\) 的全部情形。

再算首轮到下一轮的矩减少量。令
\(m=m_{0,s}(P_t,r)=|U_0^s\cap([0,P_t]\times[0,r])|\)，
\(S_{\rm cap}=P_s+Q_s\)。在目标第一排序置 \(r=Q_t-D\)，
高度由 \(P_tr/f\) 减为 \((P_tr-m)/f\)；故减少量恰为
\[
 \frac1f\int_0^{Q_t}(S_{\rm cap}-Q_t+r)m\dd r
 -\frac1{f^2}\int_0^{Q_t}(2P_trm-m^2)\dd r.
\]
因 \(m\le Ce^3/f\)，第二积分为 \(O(e^6/f^3)\)。
两个固定排序的高度参数导数为 \(O(e/f)\)，末段又有三角尾界，
因此
\[
 |U_0(P_s,Q_s)\triangle U_0(P,Q)|\le Ce^2|c|/f.
\]
帽宽改变的竖条也在此界内，对任意小 \(r\) 均成立。
代入上式的长度、核和 \(1/f\)，只花 \(Ce^4|c|/f^2\)。
外帽对减少集的裁剪因该集包含于 \(U_0\)，仍由同一长臂尾界控制。

在 \(c=0\) 时，用 Tonelli 计算减少量，核为
\[
 W_Q(Y)=P(Q-Y)+(Q^2-Y^2)/2
\]
及交换两轴的 \(W_P\)。两个排序在线性薄高积分中置薄变量为零，
误差 \(O(e^6/f^3)\)，主多项式为
\(PQ(P+Q)(5P^2+21PQ+5Q^2)/24\)。
加上 \(M_0\) 已有的 \(P^2Q^2(P+Q)/3\)，中间系数为 \(29\)，
即得 \(Q_5\) 和 \eqref{eighth:exact-cross-cap}。
实际尺度 \(P=v^3R_ep,Q=v^3R_eq,c=v^3R_ez,f=nv^2R_e\)
使其完整退款误差为
\[
 C[v^7z^2+v^8|z|+v^9]R_e^n.
\]
稳定性平方界和 Cauchy 给其全族和 \(O(v^9S)\)，
未把 RMS 偷换为逐边 \(z=O(v)\)。
合并名义两排序、三近面项和双帽退款，证明
\eqref{eighth:exact-normal-integral}。

最后作点态紧窗最大化。除以 \(n^2(1+v)^{n-4}v^6R_e^n\)，
单边净泛函是
\[
 \frac b{27}-G(b,x,y,\zeta)
 +\frac{2v}{1+v}Q_4(x,y)-\frac{2v^2}{(1+v)^2}Q_5(x,y)
 -\frac{n(n-2)v^2}{1+v}x^2y^2.
\]
横权 \(1-b\) 及根差平方有固定正罚，扰动不含这两个变量，
所以先在 \(b=1,\zeta=0\) 求最大并保留其部分罚。
固定二次强制性与扰动的 \(Cv\)-Lipschitz 界迫使最大点距基点为 \(O(v)\)。
只对这些最大点写 \(x=1/3+vh_1,y=1/3+vh_2\)，实际场不作逐点 Taylor。
由 \(Q_4(1/3,1/3)=5/243\)、其两偏导 \(10/81\)，二阶优化为
\[
 \begin{aligned}
 &-\frac{h_1^2+h_1h_2+h_2^2}{3}+\frac{20}{81}(h_1+h_2)\\
 &\quad=\frac{400}{6561}-\frac{(h_1-20/81)^2+(h_1-20/81)(h_2-20/81)
                         +(h_2-20/81)^2}{3}.
 \end{aligned}
\]
\(Q_5(1/3,1/3)=13/972\) 和三近面项分别给
\(-13/18\)、\(-n(n-2)/3\) 的相对系数，位移增益为 \(400/243\)。
将结果转成实际 \(W_N\) 时，首退款与横权差为 \(O(v)(1-b)\)，
由保留的领先罚吸收；\(\widehat b\) 截为 \(b\) 的总差更晚。
因而得到 \eqref{eighth:exact-model-price}。
首退款的因子是 \((1+v)^{-3}\)，八阶不能将其换成一。
\end{proof}

\subsubsection{整源共同高度中的全部补集}

\begin{lemma}[任意固定精度的补集结算]\label{eighth:exact-complement}
给定固定 \(\delta>0\)，可先取有限大 \(T>T_0\)，再取有限小 \(m>0\)。
令 \(\mathcal B_{T,m}\) 为正常窗外、两端深处且 \(m\le p,q\le T\)
的二轴边。在同一个 packet
\(\mathscr P=\mathcal N\cup\mathcal B_{T,m}\) 上，有固定正系数使
\begin{equation}\label{eighth:exact-complement-price}
 H+c_Dv^6D_{\mathcal N}+c_B\mathsf M_{T,m}
       +c_R\mathcal R_A^-+c_E\mathfrak E
 \le q_N(v)\mathcal C+\delta v^8S+o_\delta(v^8S).
\end{equation}
这里 \(H\) 含全部外部交对并只减一次完整 \(J\)，
\(\mathsf M_{T,m}\) 是完整非正常名义成本，\(\mathcal R_A^-\) 是
真实联合保证层之外的 raw 成本。此式对每个固定 \(n\ge2\) 成立。
\end{lemma}
\begin{proof}
正常集合、旧 \(Y_N\)、原稳定常数 \(C_*\) 和正常清洁分配先固定。
用整源 chart 写实际非负高度 \(h_A=h_N+h_B\)，名义坏高度为 \(h_B^0\)。
定义
\[
 \mathsf M=\lambda_v\sum_{\rm charts}\int h_B^0(y)\dd y,\qquad
 B_A=\mathscr B_{\mathscr P}(\mathscr D),\quad
 B_N=\mathscr B_{\mathcal N}(\mathscr D).
\]
在两个高度中使用全 packet 的同一个侵蚀 \(\epsilon_i(y)\)，得
\(B_N^-\subset B_A^-\subset D'\)。同一正部和上截断给精确常数一：
\[
 \mathcal R_N^-:=\frac va|\mathscr D\setminus B_N^-|_1
 \le\mathcal R_A^-+(1+v)^{-(n-1)}\mathsf M
 \le\mathcal R_A^-+\mathsf M.
\]
旧层精确包含于 \(B_N\)，所以
\[
 \mathcal R_A^-\le\mathcal R_N^-\le R_N+O_T(v^{10}S).
\]
不要求两类层在物理上分开。

真实非正常增量 \(\Delta_A=\lambda_v|B_A\setminus B_N|_1\) 满足
\[
 \Delta_A=\mathsf M-\mathfrak F_B-\mathfrak L_B,\quad
 \mathfrak F_B=\lambda_v\sum\int\mathcal K_B(\mathscr D;y)\dd y.
\]
纤维增量的上截断损失满足
\[
 0\le\mathfrak L_B\le
 \lambda_v\sum\int h_B^0\mathbf1_{\{H_i<h_N+h_B\}}
 \le C_Tv^3\mathsf M.
\]
最后一步仍逐名义目标证明：临界变量在 \(O_T(v^3R_i)\) 窗内，
最浅轴半宽及临界轴剩余半宽均为 \(cvR_i\) 的固定倍数；
截断只能由一个非临界自由底变量进入中分面的 \(C_Tv^4R_i\)
边带造成。该项独立于该变量，故相对损失为 \(C_Tv^3\)。
这没有裁掉源的三近面主体。

原必要性及 \(Y_N\subset B_N\) 给
\[
 \Delta_A\le(1+v)^{n-1}R_N+O_T(v^{10}S).
\]
坏退款核为 \(C_Tv^2/a\)。目标缺集按旧正常层和余量分解，
余量退款至多 \(C_TvR_N\)；旧正常层逐源为 \(O(av^5R_j^n)\)。
固定 \(p,q\ge m\) 时，两有向端完整成本由正部凸性至少为
\(c_mv^6R_e^n\)，即使根差截幅亦成立。因此
\[
 \mathfrak F_B\le C_{T,m}v\mathsf M+C_TvR_N,
\]
\[
 [1-C_{T,m}v-C_Tv^3]\mathsf M
 \le[(1+v)^{n-1}+C_Tv]R_N+O_T(v^{10}S).
\]
对每个固定 \(T,m\)，先取小 \(v\)，得
\begin{equation}\label{eighth:exact-new-window}
 \mathsf M=O_{T,m}(v^8S),\quad
 \mathfrak F_B+\mathfrak L_B=o_{T,m}(v^8S),\quad
 \limsup\frac{\mathsf M}{v^8S}\le C_*.
\end{equation}
上极限常数不随后选固定窗增长；误差常数和阈值可以增长。

由引理~\ref{eighth:exact-iteration} 的正常第二反馈，
\[
 \lambda_v|B_A|_1
 =\lambda_v|Y_{2,N}|_1+\mathsf M+o_{T,m}(v^8S).
\]
对 \(B_A^-\) 用完整强 CDF 一次，故
\[
 \sum_iJ_i\ge\lambda_v|Y_{2,N}|_1+\mathsf M
              +c_R\mathcal R_A^-+c_E\mathfrak E-o_{T,m}(v^8S).
\]
这是高度可加性和同一次实际体积余量的结果，
没有先给正常价格扣全 \(J\) 再追加另一条补集源成本。

最浅轴在坏目标 pair 内的两个排序积分为
\(\lambda_vS_i^\perp T_j^\perp P_lP_m(P_l+P_m)/(2naR_i^{n-1})\)；
在 pair 外的名义项非负，且与前者之比至多
\[
 C\frac{P_lP_m}{(P_l+P_m)w'_{iq}}\le C_Tv^2.
\]
源横宽为完整有效横宽，故旧核心的 \(O_T(v^8S)\) 误差不再出现。
整宽平均误差 \(O_T(v^{17/2}S)\)，相交根比误差为 \(Cv\mathsf M=o(v^8S)\)。
\eqref{eighth:global-coercivity} 因此保留固定
\(\gamma_\nu,c_B>0\)，在新 \(T,m\) 之前取定，并给
\[
 \sum_{\mathcal B_{T,m}}I_e-\mathsf M+c_B\mathsf M
 \le(1-\gamma_\nu)A_6W_B+o_{T,m}(v^8S).
\]

尚余近零、长角、边界和高支撑。固定
\(C_{\rm soft}\ge\max(1/m,1,12/n)\)，
\(\omega=\min(1,C_{\rm soft}\min(p,q))\)。
有界近零角满足 \(pq\le Tm\omega\)。
自适应 clean 阈值只依真实交长，不依后选软权常数，
故正常占用和 \(D_{\mathcal N}\) 保持原样，且
\[
 W_{\mathcal D},W_B\le\mathcal C-W_N+O_{C_{\rm soft}}(v^3S)
 =D_{\mathcal N}+O_{C_{\rm soft}}(v^3S).
\]
近零 gross 的归一化上极限至多 \(TmC_*\)。
长角只用固定旧正常小管及 \eqref{eighth:long-cap-fraction}：
实际目标 raw 例外给 \(CvR_N\)，面对／端移例外给
\(Cv(\mathcal F+\mathfrak E)\)，其余实际长短条避开旧正常层，
固定重数给
\[
 \sum_{\mathcal L_T}I_e\le(C_L/T)R_N+
                 Cv(R_N+\mathcal F+\mathfrak E).
\]
\(C_L\) 与新 \(T,m\) 无关；不要求长条避开新整源大管，
也不从本次 \(J\) 再购买这些已稳定的几何数值界。
边界 epsilon 分解如引理~\ref{eighth:exact-high-gross}，
其二轴 gross 为 \(O(v^9S)\)。高支撑为 \(o(v^8S)\)。

二维须单列实际桥：唯一轴对是 \(\{1,2\}\)，最浅轴必在 pair 内，
没有非临界自由底变量，所以 \(\mathfrak L_B=0\)，pair 外名义项为空。
全部原始横积为一，两个不等根有向排序、真实 clip 和上面小算子仍在。
其容量由实际四角／四胞及根桥给
\(\mathcal C=\sum_iR_i^2,\ W_N=\mathcal C-D_{\mathcal N}+O(v^3S)\)，
不跨维引用高维 helper。原长短条至少长轴深度超过正常小管，
同一二维原见证的轴帽矛盾、截宽和固定重数证明给同一 \(C_L/T\) 尾；
高支撑为空，边界二轴仍为九阶。因此以上补集结算也在二维原位成立。

给 \(\delta\) 后先取 \(T\) 使 \(C_LC_*/T\) 足够小，
再取 \(m\) 使 \(TmC_*\) 足够小，随后固定 \(C_{\rm soft}\)，最后 \(v\to0\)。
合并得到
\[
 H+\text{所列正量}\le
 \mathsf P_N+(1-\gamma_\nu)A_6W_B+\delta v^8S+o_\delta(v^8S).
\]
用 \eqref{eighth:exact-model-price} 和槽恒等式，右侧至多
\[
 q_N\mathcal C-
 [q_N-(1-\gamma_\nu)A_6]D_{\mathcal N}
 +\delta v^8S+o_\delta(v^8S).
\]
括号至少为固定 \(c_Dv^6\)，即得结论。
槽欠量只保留一次，没有另加正常模型单独转换所得的第二份槽收入。
\end{proof}

\subsubsection{二维小半径源的常数一回收}

二维尚有同阶的盒内小源费用。以下局部引理处理原见证跨盒边界的情形。

\begin{lemma}[角窗中的精确局部费用]\label{eighth:exact-small-box}
设矩形 \(\mathcal R\) 各边介于 \(cL,CL\)，
\(T\subset\mathcal R\setminus E\)、\(Z\subset\mathcal R\) 可测。
\(Z\) 位于四角边长 \(\epsilon L\) 的块内，且
\[
 D=\mathcal R\setminus T,\quad z=|Z\setminus T|,\quad
 d=|D\setminus Z|,\qquad |Z|,d\le bL^2.
\]
若 \(b,\epsilon\to0\)、\(\sqrt b/\epsilon\to0\)，
任意可测的 \(\mathcal S\subset T\) 的原严格见证半径取有限正值并满足
\(r(x)\le AL\) 时，
\[
 |\mathcal S|\le\frac a{1-a}z+C_{c,C,A}ad.
\]
价格常数不依最小原半径或有限列表长度。
\end{lemma}
\begin{proof}
从整个 \(\mathcal S\) 而非有限 probe 中按半径值降序选尚未被此前闭
\(Q\) 覆盖的中心。新中心与旧中心距大于两个半径最大值。
有界性和本次正最小半径的 packing 保证有限终止，因此原 \(Q\)
覆盖全部 \(\mathcal S\)。选中数可依最小半径，价格不使用这个数界。
正式引理~\ref{lem:finitegreedy} 已使全局原函数为有限正值 Borel 函数；
限制到这里的可测集合不改变半径，所以没有对一般可测函数假设正下界。

严格见证给 \(|Q(x,r)\cap T|<ar^2\)。
若 \(r\ge\ell_{\min}\)，盒内交面积至少 \(\ell_{\min}^2\)，
而 \(|D|\le2bL^2\)、\(r\le AL\)，小参数时矛盾。
所以 \(r<\ell_{\min}\)，每轴盒内交长至少 \(r\)，从而
\[
 (1-a)r^2<|D|\le2bL^2,\qquad r=o(\epsilon L).
\]
这是每一点的原半径结论，未要求原 \(Q\) 全在盒内。

将四角窗扩大到边长 \(2\epsilon L\)，并为 \(V\)。
每角分别作上述全点选择，对每个选中中心取向盒内的
\(P=x+\sigma[0,r]^2\)。它在盒内且仍落于该角 \(3\epsilon L\) 块。
同符号两 \(P\) 不交：实现强中心分离的某轴中心距大于两个半径，
故该轴的同向区间不交。不同角的 \(3\epsilon L\) 块也分离。
所以近角的所有 \(P\) 物理不交，且
\(|P\cap D|\ge(1-a)r^2\)。用原 \(Q\) 的全点覆盖得
\[
 |\mathcal S\cap V|\le\sum|Q\cap T|<a\sum r^2
 \le\frac a{1-a}|D|=\frac a{1-a}(z+d).
\]
远角源点到 \(Z\) 至少 \(\epsilon L\)，其原 \(Q\) 与 \(Z\) 不交。
按两个坐标位于盒中点哪侧分四组，每组取朝中心的固定符号。
半个盒宽足以容纳 \(r\)，故对应 \(P\subset\mathcal R\setminus Z\)；
每组全点选择后 \(P\) 仍不交，四组仅向 residual 收费：
\[
 |\mathcal S\setminus V|\le4a\,d/(1-a).
\]
合并给 \(az/(1-a)+5ad/(1-a)\)。正常 \(z\) 没有四倍重数。
有限选择不要求局部集合紧，所有集合均可测。
\end{proof}

\begin{lemma}[实际删除量及基线导数]\label{eighth:exact-small-source}
在同一二维第八层，
\begin{equation}\label{eighth:exact-inside}
 (1+O(v))\frac va\sum_{\rm deep}\delta_{i,\rm in}
 \le(2K_2+O(v))v^8S,\qquad K_2=4/27.
\end{equation}
边界、环境及深处盒外费用分别为 \(O(v^{10}S),O(v^{10}S),O(v^9S)\)。
实际基线中的删除系数可单独保留为 \(1+O(v)\)，
原重复量正预算及源缺量望远镜不变。
\end{lemma}
\begin{proof}
取原入向面邻构成的扩展盒 \(\mathcal R_i\)，仍各边可比于 \(vr_i\)，
并记
\[
 q_i=|\mathscr D_i\setminus Y_{i,N}|+|\mathcal R_i\setminus(x_i+W_i^0)|.
\]
入向扩张体积及稳定预算给
\(\sum q_i\le(a/v)R_N+Cv\mathcal F=O(av^7S)\)。
坏盒 \(q_i>av^5r_i^2\) 的原权为 \(O(v^2S)\)。
对它们只用原 dyadic 引理和 \(F_i\subset T_i\)，其 inside 总量为
\[
 Ca\sum_{\rm bad}(|Y_{i,N}\setminus T_i|+q_i)=O(a^2v^7S).
\]
该项乘 \(v/a\) 后为十阶，不污染领先系数。

好盒上 \(|\mathscr D_i|\le Cav^5r_i^2\)，
\(\varepsilon_i\mu_i\le|\mathscr D_i|\) 给 \(\varepsilon_i\le Cv^5\)。
中点正盈余给端移 \(\Delta_i\le C\varepsilon_i vr_i\le Cv^6r_i\)；
扩张面积除以原盒另一轴宽给每侧扩张长度也为 \(O(v^6r_i)\)。
所以旧正常层确在扩展盒的四个 \(C_Nv^3r_i\) 角窗内。
取上引理的 \(L=vr_i,b=Cv^5,\epsilon=C_Nv^2\)，
\(\sqrt b/\epsilon=O(v^{1/2})\to0\)，并有
\(|(\mathcal R_i\setminus T_i)\setminus Y_{i,N}|\le q_i\)。
原小半径仍小于 \(Avr_i\)，引理还逐点推出其为 \(O(v^{7/2}r_i)\)。
故
\[
 \sum_{\rm deep}\delta_{i,\rm in}
 \le\frac a{1-a}\sum_i|Y_{i,N}\setminus T_i|+Ca^2v^7S.
\]

正常层体积的领先常数只需名义两端积分：
二维有向两排序面积为 \(PQ(P+Q)/(4aR_i)\)。
在 \(p=q=2/[3(1+v)],z=0\)，无向两端除 \(av^5R_e^2\)
为 \(8/[27(1+v)^3]\)。稳定性给加权参数平方差 \(O(v^2S)\)，
Cauchy 给一阶差 \(O(vS)\)；紧窗积分 Lipschitz，因而
\[
 \sum_i|Y_{i,N}|\le2K_2av^5\sum_{\mathcal N}R_e^2+O(av^6S)
 \le(2K_2+O(v))av^5S.
\]
最后用二维实际四槽 \(\mathcal C=\sum R_i^2\le S\)。
乘 \(v/a\) 即得 \eqref{eighth:exact-inside} 的体积部分；
未借正常八阶积分。

稳定性保留的 \(P_I+P_\partial=O(v^8S)\) 给精确边界小源式
\((v/a)\sum_{\partial I}\delta_i\le CaP_\partial=O(v^{10}S)\)；
环境 \(Ca(\mathcal F+P_\partial)\) 同阶，盒外为 \(Cv\mathcal F=O(v^9S)\)。

最后必须核基线系数。令 \(m_i^0=|T_i|\)、\(\delta_i=m_i^0-\mu_i\)、
\(f_i(m)=(R_i+\sqrt m)^2\)，实际质量给
\((1-Av)g_i\le\mu_i\le m_i^0\le g_i=aR_i^2\)。
在整个真实区间，
\[
 \frac{f_i'(m)}{v/a}=v+\frac{vR_i}{\sqrt m}
 \le v+(1-Av)^{-1/2}\le1+(1+2A)v.
\]
故一次真实积分可拆为
\[
 f_i(\mu_i)\ge A_vg_i/a
 -(1+O(v))(v/a)(g_i-m_i^0)
 -(1+O(v))(v/a)\delta_i.
\]
\(\sum(g_i-m_i^0)=\mathcal D_h+\mathfrak d_h^0\) 仍精确；
环境扣一次 \(\mathcal D_h\) 后其正系数至少 \(1/(2a)\)。
只对 \(\mathfrak d_h^0\) 使用原固定 \(b\) 及
\eqref{eighth:depth}，删除量单列不被 \(b\) 放大。
\(|X_i|\ge f_i(\mu_i)+J_i\) 与 Bonferroni 仍给
\(\sum f_i(\mu_i)-H_h\)，所以 \(J\) 和删除质量各只计一次。
\end{proof}

\subsubsection{八层中心化与真实分母}

\begin{proposition}[八阶精确上界]\label{eighth:upper}
对每个固定 \(n\ge2\)，
\[
 \limsup_{t\downarrow0}
       \frac{C_n(1-t)-P_{7,n}(t)}{av^8}\le\widehat J_n .
\]
\end{proposition}
\begin{proof}
先沿任意 \eqref{eighth:exact-competitive} 序列工作。
第八层取引理~\ref{eighth:exact-complement} 的同一个正常模型，
不再加入正常价格单独结算的第二份 \(D_{\mathcal N}\)。
完整槽乘积和质量恒等式给
\[
 \mathcal C=B_n(1+v)^{n-2}\sum_iR_i^n+O(v^{7/2}S),\qquad
 \sum_iR_i^n=(M'+\mathcal D_8)/a
                 \le\mathcal V/A_v+O(v^6S).
\]
乘 \(q_N=O(v^6)\) 后误差晚于八阶，因此令
\[
 Q_{8,N}(v)=K_n\left[
 \frac{v^6}{(1+v)^4}+\frac{10v^7}{9(1+v)^5}
                         +\frac{c_{8,n}v^8}{(1+v)^2}\right].
\]
首退款的分母五次幂不能改成二次幂。

对一般维数，写 \(m/g_i=u\in[1-A_8v,1]\)，
\[
 \frac{d(R_i+m^{1/n})^n/dm}{v/a}
                   =(v+u^{-1/n})^{n-1}=1+O_{n,A_8}(v).
\]
将 \(g_i-\mu_i=(g_i-m_i^0)+\delta_i\) 的两段真实积分分开，
原基线仍有重复量净正系数至少 \(1/(2a)\)。
只对未删除源缺量 \(\mathfrak d_h^0\) 采用固定 \(b\)，
删除量保留 \(1+O(v)\) 系数。
当 \(n\ge3\)，原 dyadic 内盒、扩张及环境绝对项
\(v^{n+6}S=o(v^8S)\)，盒外 \(Cv\mathcal F=O(v^9S)\)。
当 \(n=2\)，用引理~\ref{eighth:exact-small-source}。
令
\[
 \kappa_n=\frac8{27}\mathbf1_{\{n=2\}},\qquad
 Q_8(v)=Q_{8,N}(v)+\kappa_nv^8 .
\]
\(S\le\mathcal V/(1-a)\) 把唯一 inside 费用换成
\(\kappa_nv^8\mathcal V+o(v^8S)\)。这是独立的删除成本，
不是下面真实分母产生的二维项。
于是恢复后的第八层可选并集误差上包
\[
 E_8=Q_8(v)\mathcal V+\delta v^8S+o_\delta(v^8S).
\]
若实际误差更小，只选此上包；没有将实际容量不等式写成等式。

第七前缀在本次八层族重放粗价格，使用本层实际对象而非第八层重命名，
可取 \(E_7=Q(v)\mathcal V+O(v^8S)\)。
这里 \(Q-Q_8=O(v^8)\)，故 \(E_7-Q_8\mathcal V=O(v^8S)\)。
第六前缀仍取领先上包 \(E_6=(K_n+o(1))v^6S\)，其中心化差为
\(o(v^6S)\)；前五层只取 \(Cv^6S\)。
所有这些前缀与原见证、\(I\) 和统一 CDF 在开始即固定。

回到公共账本 \eqref{eighth:ledger}，仍用
\(\tau=bv/c_D'\)、\(\Omega_8=\sum_{j=0}^7\tau^j\)。
\eqref{eighth:telescope} 精确消去中间源缺量，初端除 \(a\Omega_8\)
为 \(O(v^9S)\)。必须对整个权平均减同一个中心价：
\[
 \bar E-Q_8\mathcal V=
 \frac{[E_8-Q_8\mathcal V]+\tau[E_7-Q_8\mathcal V]
 +\tau^2[E_6-Q_8\mathcal V]
 +\sum_{h=1}^5\tau^{8-h}[E_h-Q_8\mathcal V]}{\Omega_8}.
\]
四类贡献依次为
\(\delta v^8S+o_\delta(v^8S),O(v^9S),o(v^8S),O(v^9S)\)。
图条件吸收固定边界和内部体积价，剩余正回收和 \(f_8\) 可丢，得到
\begin{equation}\label{eighth:exact-common-ledger}
 \mathcal V\ge A_vM/a-Q_8(v)\mathcal V
                   -\delta v^8S-o_\delta(v^8S).
\end{equation}
外部族为空时所有源模型和交量取零，不定义空 CDF；
内部共享质量及突出回收直接给同方向。

令 \(x=M/V\)，利用 \(S/V\le(1-a)^{-1}\)，直接移实际 \(M\) 得
\[
 1+x\le\frac{A_v}{A_v-a(1+Q_8)}
                +C_n\delta av^8+o_\delta(av^8).
\]
分母小参数时统一远离零，与 \(L_n\) 的精确差是
\[
 \frac{aA_vQ_8}{(A_v-a)(A_v-a-aQ_8)}.
\]
\(Q_8=O(v^6)\)，再除 \(a\) 后将末分母换成 \(A_v-a\)
只损 \(O(aQ_8^2)=o(v^8)\)。所以只须展开
\[
 Q_8\frac{A_v}{(A_v-a)^2},\qquad
 Q_{8,N}=K_n[v^6-\tfrac{26}{9}v^7+(40/9+c_{8,n})v^8]+O(v^9).
\]
真实分母满足
\[
 \frac{A_v}{(A_v-a)^2}
 =1-nv+\left[\frac{n(n+1)}2+2\mathbf1_{\{n=2\}}\right]v^2+o(v^2).
\]
一般整数维数的余项可令 \(a=v^2w,\ 0\le w\le1\)：
分母在共同邻域远离零，三阶 Taylor 余项对 \(w\) 一致；
随后 \(n\ge3\) 代 \(w=v^{n-2}\to0\)，二维代 \(w=1\)。
故八阶系数为
\[
 K_n\left[\frac{40}{9}+c_{8,n}+\frac{26n}{9}
                  +\frac{n(n+1)}2+2\mathbf1_{\{n=2\}}\right]+\kappa_n
 =\widehat J_n.
\]
其中二维分母已贡献 \(2K_2=8/27\)，inside 另贡献
\(\kappa_2=8/27\)，两者来源不同。

最后处理全部输入上确界。若上极限失败，取固定 \(\eta>0\) 及
\(v_j\downarrow0\)，按上确界取实际输入，使归一化余量超过
\(\widehat J_n+\eta\)，再以 \(o(Vav^8)\) 精度取外部严格 halo 紧集。
密度定理保证输入几乎处处在自身 halo，外部 halo 有限可测，内正则性适用。
该序列满足 \eqref{eighth:exact-competitive}，例如可取固定 \(B=1\)。
先固定这条序列和原稳定构造，再取固定 \(\delta\) 使
\(C_n\delta<\eta/4\)，依次取有限 \(T,m\)，最后 \(v\to0\)，即矛盾。
逐序列小量模量可以不同，未将未证统一速率移出上确界。
\end{proof}

\subsubsection{对称斜格孔与两次角反馈}

\begin{proposition}[宏观孔的有限输入构造]\label{eighth:macro-lower}
对每个充分小的 \(t>0\)，存在有界有限正测度输入 \(E_t\)，使
\begin{equation}\label{eighth:finite-lower}
 \frac{|\mathcal H_n(1-t;E_t)|}{|E_t|}
 =P_{7,n}(t)+J_nav^8+O_n(av^9).
\end{equation}
\end{proposition}
\begin{proof}
置 \(u=t^{1/n}\)、\(g=(1+2u)^{-1}\)，取
\[
 e=\kappa u^3,\quad T=(1-e)I+eJ,\quad d=\det T,\quad
 w=ug+\Lambda u^6,\quad \rho=1-2w,\quad M=a\rho^n ,
\]
其中 \(J\) 为全一矩阵。先把 \(\kappa\) 限于一个固定正紧区间、
\(\Lambda\) 限于固定有界区间，下面所有误差在该区间上一致；
最后给出参数并先选固定安全常数、再取小 \(u\)。
在 \(C_w=[-w,w]^n\) 上，对每个 \(x\) 用同一规则
\[
 \begin{aligned}
 &s=w-\|x\|_\infty,\quad r=\rho+s,\quad\\
 &F(s)=\frac{a}{(2w)^{n-2}}[(\rho+s)^n-\rho^n],\quad\\
 &f=F'(0),\quad h=e^2/f,\quad \delta=h/e.
 \end{aligned}
\]
有 \(f=4nu^2g[1+O_n(u^5)]\)、\(h=O_n(u^4)\)、\(\delta=O_n(u)\)。

先完整排除其他邻孔。因 \(r+w=1-\|x\|_\infty\)、
\(\|Tk\|_\infty\ge(1-(n-1)e)\|k\|_\infty\)，
小 \(u\) 时外包孔分离且 \(\|k\|_\infty\ge2\) 的孔不能进入。
对 \(k\in\{-1,0,1\}^n\)，若 \(\sum k_l\ne0\)，取与和同号的非零坐标，
其格中心模为 \(1+e(|\sum k_l|-1)\ge1\)，至多边界接触。
只有非零平衡指标 \(\sum k_l=0\) 进入，支撑大小 \(j\) 为偶数，
\(Tk=(1-e)k\)，各支撑交长至多 \(e\)，其余取整宽。
又 \(r\ge\rho>2w\)，故自身全孔始终包含。

对任何 Borel 删除集 \(U\subset C_w\)，在整个外包孔定义
\[
 \begin{aligned}
 R_U(x)&=\frac1{(2w)^{n-2}}\sum_{k\ne0}
       |(Tk+(C_w\setminus U))\cap Q(x,\rho+s)|,\quad\\
 \mathcal B(U)&=\{x\in C_w:s=0\text{ 或 }R_U(x)\ge F(s)\}.
 \end{aligned}
\]
有限邻孔清单和方体平移、变半径的对称差连续性保证这些坏集为 Borel。
映射 \(\mathcal B\) 反单调，置
\begin{equation}\label{eighth:three-sets}
 U_0=\mathcal B(\varnothing),\qquad U_1=\mathcal B(U_0),\qquad
 U_2=\mathcal B(U_1).
\end{equation}
则 \(U_1\subset U_2\subset U_0\)，从而 \(\mathcal B(U_2)\subset U_2\)。
所以只要 \((2w)^n-M>|U_2|\)，任何在 \(C_w\setminus U_2\)
中选出的质量 \(M\) 的孔都具有严格见证。
每次容量仍用同一个实际自身质量 \(M\)，无需求未知孔形的固定点。

所有归一化外包项至多 \(e^2\)，数目少于 \(3^n\)，故三坏集内点的
\(s\le h_H:=F^{-1}(3^ne^2)\le3^nh\)。
忽略零测并列面，选全局最浅轴 \(i\)，在固定符号区内写
\(d_l=w-|x_l|=s+y_l\)、\(y_i=0\)。
雅可比为一。活动支撑只能朝相同源面符号，穿透为 \((e-y_l)_+\)，
目标交集相对于面对端的区间不随 \(s\) 变；
非支撑方向全宽进入。因此可将它记为 \(R_U(y)\)，
实际坏层厚度是
\[
 \min\{F^{-1}(R_U(y)),\,w-\max_{l\ne i}y_l\}.
\]
计算 \(|U_2|/a\) 到 \(u^7\) 时，可统一替换为 \(R_{U_1}(y)/f\)，
绝对误差 \(O_n(u^8)\)。具体地，\(F'(s)/f=1+O_n(u^4)\)，
外包支撑 \(j\) 的积分在最浅轴位于支撑内、外时，分别至多
\[
 O_n\!\left(\frac{(2w)^{2-j}e^{2j-1}w^{n-j}}{af}\right)=O_n(u^{4j-3}),
 \quad
 O_n\!\left(\frac{(2w)^{2-j}e^{2j}w^{n-1-j}}{af}\right)=O_n(u^{4j-1}).
\]
总量为 \(O_n(u^5)\)，\(j\ge4\) 仅从 \(u^{13}\) 起。
若 \(y_l>w-h_H\) 触发厚度截断，该轴必在所计支撑外，
否则 \(y_l>e\) 使交量零。逐项外包积分在该轴恒定，
截去的边带只损相对 \(O_n(h_H/w)=O_n(u^3)\)，给绝对 \(O_n(u^8)\)。
线性化后对所有同时进入的目标求和，不将坏体积按不交角集作错误加法。

下面计算两支撑目标的第二轮反馈。
在固定短角以 \(X,Y\) 表示两个物理面深除以 \(e\)，取 \([0,1]^2\)。
记 \(\overline m_l(q)\) 为 \(U_l\) 在该目标帽
\([0,e]\times[0,eq]\) 乘全部其他宽中的质量，除以
\(e^2(2w)^{n-2}\)。
轴置换和全反射保持三个坏集，所以交换排序与相反短角同价。
令 \(D=|X-Y|\)、\(q=1-D\)，定义平面模型
\[
 P_0=\{\min(X,Y)\le\delta q\},\quad
 m_0(q)=|P_0\cap([0,1]\times[0,q])|,\quad
 P_1=\{\min(X,Y)\le\delta(q-m_0(q))\}.
\]
有 \(0\le m_0(q)\le q\)。
若其余坐标距各面均 \(>2e\)，只有该短角可能进入，
\(U_0\) 的截线为
\(F^{-1}(e^2q)/e=\delta q+O_n(\delta u^4)\)。
非中央切片须另计：某横向面深至多 \(h_H\) 的比例为
\(O_n(h_H/w)=O_n(u^3)\)；
最浅横向面深在 \((h_H,2e]\) 的比例为 \(O_n(u^2)\)，
但坏点还必须有 \(\min(X,Y)\le h_H/e=O_n(u)\)，故仍为 \(O_n(u^3)\)。
平面模型同样受此面积界控制，因而
\(\sup_q|\overline m_0(q)-m_0(q)|=O_n(u^3)\)。
中央 \(U_1\) 截线为
\(F^{-1}(e^2[q-\overline m_0(q)])/e\)，
与 \(P_1\) 差 \(O_n(\delta u^3)\)；
非中央因 \(U_1\subset U_0\) 仍有同样总误差。二维横向例外为空。

于是 Tonelli 给
\begin{equation}\label{eighth:second-moment}
 2\int_0^1\overline m_1(q)\dd q
 =I_1+O_n(u^3),\qquad
 I_1=\int_{P_1}(2-X-Y)\dd X\dd Y
 =\frac53\delta-\frac{13}{4}\delta^2+O_n(\delta^3).
\end{equation}
核对最后系数：在 \(X\le Y\) 的排序，\(Y=X+D\)、\(0\le X\le q\)，
\(P_0\) 取 \(X\le\delta q\)，\(P_1\) 取 \(X\le\delta(q-m_0(q))\)。
所以
\[
 I_0=\frac53\delta-\frac23\delta^2,\qquad
 I_0-I_1=2\delta\int_0^1(1+q)m_0(q)\dd q+O(\delta^3).
\]
再以 Tonelli 写后一个积分为 \(\int_{P_0}W(Y)\)，
\(W(Y)=3/2-Y-Y^2/2\)。
两个排序分别将 \(W(Y)\) 换为 \(W(D)\)、\(W(0)\)，
Lipschitz 误差为 \(O(\delta^2)\)，主项为
\[
 \delta\int_0^1(1-D)[W(D)+W(0)]\dd D
 =\frac{31}{24}\delta .
\]
故 \(-2/3-31/12=-13/4\)，重合臂和排序到面距的区别均已包含。

最浅轴在某两支撑内时，对全部轴对的两个相反短角和两排序求和，贡献为
\[
 P\left[1-\frac53\delta+\frac{13}{4}\delta^2\right]+O_n(u^8),
 \qquad P=\frac{2B_n(2w)^{n-2}e^3}{af}.
\]
最浅轴不在该对内时，两目标穿透的面积积分为 \(e^4/4\)。
整个目标面窗中 \(U_1\subset U_0\) 的相对删量至多
\(C_n(h_H/e+h_H/w)=O_n(u)\)，所以此项换成完整面窗只损 \(O_n(u^8)\)。
计入 \(n-2\) 种最浅轴及两面符号，得到必须保留的三面项
\[
 T_3=\frac{B_n(n-2)(2w)^{n-3}e^4}{af},
 \qquad
 \frac{|U_2|}{a}
 =P\left[1-\frac53\delta+\frac{13}{4}\delta^2\right]+T_3+O_n(u^8).
\]
二维 \(T_3=0\)。更高平衡支撑已由前述积分幂界排除到误差内。

统一展开为
\[
 \begin{aligned}
 P&=\frac{B_n\kappa^3}{8n}u^5g^{n-3}[1+O_n(u^5)],\quad\\
 \delta&=\frac{\kappa}{4ng}u[1+O_n(u^5)],\quad\\
 T_3&=\frac{B_n(n-2)\kappa^4}{32n}u^7g^{n-4}[1+O_n(u^5)].
 \end{aligned}
\]
实际可删量为
\(((2w)^n-M)/a=n\Lambda u^5g^{n-2}+O_n(u^{10})\)。
因此取
\begin{equation}\label{eighth:lower-mass}
 \begin{aligned}
 \Lambda={}&\frac{B_n\kappa^3}{8n^2g}
 -\frac{5B_n\kappa^4}{96n^3g^2}u\\
 &+\left[\frac{13B_n\kappa^5}{512n^4g^3}
        +\frac{B_n(n-2)\kappa^4}{32n^2g^2}\right]u^2+c_nu^3.
 \end{aligned}
\end{equation}
前三项逐项抵消坏量预算；余量为
\(u^8[nc_ng^{n-2}+O_n(1)]\)。
先在预定紧区间统一误差，再选固定大整数 \(c_n\)，最后小 \(u\)，
有 \((2w)^n-M>|U_2|\)。
沿第一坐标连续分位切片选 Borel 孔
\(P_u\subset C_w\setminus U_2\)、\(|P_u|=M\)。
因 \(\mathcal B(U_2)\subset U_2\)，每个孔点的原规则半径严格满足
\[
 \left|\bigcup_k(Tk+P_u)\cap Q(x,\rho+s)\right|
 \le M+(2w)^{n-2}R_{U_2}(x)<a(\rho+s)^n.
\]

剪切参数最后取
\begin{equation}\label{eighth:lower-shear}
 \kappa=\frac{8n}{3}g^2+\frac{320n}{81}g^3u .
\end{equation}
它留在前述固定正紧区间；\(\Lambda\) 前三项有界，
选安全常数后只须再令 \(c_nu^3\le1\)，故无选参循环。
用 \(m_*=ag^n\)、\(d=1-B_n\kappa^2u^6+O_n(u^9)\)，精确相减
\[
 \frac d{d-M}-L_n(t)=\frac{M-dm_*}{(d-M)(1-m_*)}.
\]
分子除以 \(au^6\) 为
\[
 B_ng^n\left[
 \kappa^2-\frac{\kappa^3}{4ng^2}
 +\frac{5\kappa^4}{48n^2g^3}u
 -\left(\frac{13\kappa^5}{256n^3g^4}
       +\frac{(n-2)\kappa^4}{16ng^3}\right)u^2\right]+O_n(u^3).
\]
零阶在 \(\kappa_*=8ng^2/3\) 最优，二阶导为 \(-2B_ng^n\)。
一阶导除以 \(2B_ng^n\) 在该点为 \(320ng^3/81\)，故
\eqref{eighth:lower-shear} 的位移增加 \((400/243)K_nav^8\)。
逆分母在二维是 \(1+8u^2+O(u^3)\)，在 \(n\ge3\) 为 \(1+O_n(u^3)\)。
再展开 \(g\)，所得八阶系数满足
\[
 \frac{J_n}{K_n}
 =\frac{(n+4)(n+5)}2-\frac{10}{9}(n+5)-\frac{13}{18}
       -\frac{n(n-2)}3+\frac{400}{243}+2\mathbf1_{\{n=2\}},
\]
等于定理所写正多项式；尤其 \(J_2=15694/6561\)。

最后取整数 \(N=\lceil u^{-10}\rceil\ge4\)、
\(\Omega_N=T([-N,N]^n)\)、
\(E_t=\Omega_N\setminus\bigcup_k(Tk+P_u)\)，保留全部边界切孔。
周期积分精确给 \(|E_t|=(2N)^n(d-M)>0\)。
\(w+r\le1\)、\(\|T^{-1}\|_{\infty\to\infty}<2\) 保证
\(\|k\|_\infty\le N-3\) 的孔及其严格见证均在 \(\Omega_N\) 内。
输入密度点也进入严格 halo；
凸外盒外任一中心的方体至多半体积在分离半空间内，所以小 \(t\) 下
外部点不进入。有限比值因而为
\(d/(d-M)+O_n(a/N)\)。
安全余量首次影响 \(av^9\)，截断误差更小，
遂得对每个充分小实数 \(t\) 的 \eqref{eighth:finite-lower}。
\end{proof}

\subsubsection{二维两尺度孔与匹配下界}

\begin{proposition}[精确八阶的有限下界]\label{eighth:lower}
对每个充分小的 \(t>0\)，存在有界正有限测度输入 \(E_t\)，使
\[
 \frac{|\mathcal H_n(1-t;E_t)|}{|E_t|}
       =P_{7,n}(t)+\widehat J_nav^8+O_n(av^9).
\]
\end{proposition}
\begin{proof}
\(n\ge3\) 已由命题~\ref{eighth:macro-lower} 给出。下面固定二维，
沿用其 \(u,g,\kappa,e,T,d,C_w,\rho,M,F\) 和真实坏集算子
\(\mathcal B\)，始终 \(M=a\rho^2\)，不因加入微孔而重算根。
记
\[
 \Lambda_0=\frac{\kappa^3}{32g},\qquad
 \Lambda_1=-\frac{5\kappa^4}{768g^2},\qquad
 \Lambda_2=\frac{13\kappa^5}{8192g^3},
\]
并在宏孔参数中改取
\begin{equation}\label{eighth:exact-micro-budget}
 \Lambda=\Lambda_0+\Lambda_1u+\Lambda_2u^2-a\Lambda_0+cu^3.
\end{equation}
\(\kappa\) 仍取 \eqref{eighth:lower-shear} 的二维值。
前面的全部实际集合估计对有界 \(\Lambda\) 一致，
因此此改动后可重新明确构造 \(U_0,U_1,U_2\)。

二维只有 \(k=(1,-1),(-1,1)\) 的外孔可产生正交量，
浅带中四个最浅轴 chart 的目标帽为
\([0,e]\times[0,(e-y)_+]\) 或交换两轴，独立于 \(s\)。
坏层高度是连续函数
\(\min(F^{-1}(R(y)),w-y)\)；
帽与任意固定 Borel 集的交量连续，故 \(U_2\) 闭且边界零测。
其体积等于内部体积，并有前面已证的统一展开
\[
 |U_2|/a=2\Lambda_0u^5+2\Lambda_1u^6+2\Lambda_2u^7+O(u^8).
\]
该首项正，允许在真实内部放入微孔，而非在形式极限轮廓中放孔。

固定当前小 \(u\)，取 \(t'=t(1-u)<t\)、\(v'=2\sqrt{t'}\)、\(a'=4t'\)。
对 dyadic 间距 \(h\)，只保留
\(q\in h\mathbb Z^2\) 中满足 \(Q(q,h)\subset\operatorname{int}U_2\) 的点，
并取微孔半宽和并集
\[
 w_\mu=\frac{hv'}{2(1+v')},\qquad
 H_\mu=\bigcup_q(q+[-w_\mu,w_\mu]^2).
\]
这里 host \(Q(q,h)\) 可以重叠；体积计数使用不交基本胞
\(q+[-h/2,h/2)^2\)。当 \(h\to0\)，有效基本胞总体积趋于 \(|U_2|\)：
距内部补集大于 \(3h/2\) 的点都被有效胞覆盖，内部边带趋零。
故可取第一个满足 \(h\le u^{20}\) 且有效胞体积至少
\((1-u)|U_2|\) 的 dyadic 间距，选择有限终止。
微孔不交，并有
\begin{equation}\label{eighth:exact-micro-mass}
 |H_\mu|=a|U_2|[1+O(u)].
\end{equation}

每个微孔的所有点 \(z\) 保留短半径
\[
 r_\mu(z)=h-w_\mu-\|z-q\|_\infty.
\]
小 \(u\) 时 \(r_\mu\ge h-2w_\mu\ge2w_\mu>0\)，故含自身全微孔。
其他微格孔在某轴的距离至少 \(r_\mu\)，至多边界相碰；
且 \(\|z-q\|_\infty+r_\mu=h-w_\mu<h\)，所以整个见证在 host 内。
其缺陷率至多
\[
 \left(\frac{w_\mu}{h-2w_\mu}\right)^2=t'<t.
\]
这是原阈值的逐点严格见证，含微孔边界点。

必须保护宏观见证所见的微孔质量。固定已经构造的
\(H_\mu\subset U_2\)，定义
\[
 \begin{aligned}
 \mathcal B_\mu(V)=\{s=0\}&\cup\biggl\{x\in C_w:\\
 &\sum_{k\ne0}|(Tk+(C_w\setminus V))\cap Q(x,\rho+s)|\\
 &\quad+\sum_{k\ne0}|(Tk+H_\mu)\cap Q(x,\rho+s)|\ge F(s)\biggr\}.
 \end{aligned}
\]
下面以 \(R_{H_\mu}(x)\) 记此第二个其他周期微孔质量和，\(R_U\) 沿用宏观定义。
置 \(U=U_2\cup\mathcal B_\mu(U_2)\)。
因 \((C_w\setminus U_2)\cup H_\mu\subset C_w\) 且两部分不交，
\(\mathcal B_\mu(U_2)\subset U_0\)；
反单调性给 \(\mathcal B_\mu(U)\subset\mathcal B_\mu(U_2)\subset U\)。
这是一遍显式集合运算，没有未知孔形固定点。

在上述四个 chart 中，新旧目标帽质量差为
\(|H_\mu\cap\Pi(y)|\)，只有 \(0\le y\le e\) 可能非零。
\(F'\ge f=2a\rho\)，逆函数和真实纤维上截断为 \(1/f\)-Lipschitz。
利用 \(\mathcal B(U_2)\subset U_2\)，得
\[
 |U|\le|U_2|+\frac1f\sum_{\text{四图}}\int_0^e
                         |H_\mu\cap\Pi(y)|\dd y
 \le|U_2|+\frac{4e}{f}|H_\mu|
 \le|U_2|+Cu|H_\mu|.
\]
无需假设每个宏观帽中的微孔恰有均匀密度。
深源的完整外包交量仍不足 \(F(s)\)，故四图外没有漏掉新的坏点。

记可删总量 \(D=(2w)^2-M\)，其精确式为
\[
 D/a=2\Lambda u^5+\Lambda^2u^{10}(1-4u^2).
\]
由 \eqref{eighth:exact-micro-budget}、
\eqref{eighth:exact-micro-mass}，
\[
 [D+|H_\mu|-|U|]/a\ge2cu^8-Cu^8.
\]
先在固定参数紧区间取统一 \(C\)，再选大固定 \(c\)，最后小 \(u\)
使 \(cu^3\le1\)，所以选参没有循环且余量严格正。
\(M-|H_\mu|>0\)，可在开集 \(C_w\setminus U\) 内先取质量更大的紧集，
再用第一坐标连续分位切出恰为 \(M-|H_\mu|\) 的紧宏孔 \(A_\mu\)。
总孔 \(P_\mu=A_\mu\cup H_\mu\) 的实际质量精确为 \(M\)。

宏孔 \(x\in A_\mu\) 的原长半径 \(\rho+s\) 含整个本胞 \(P_\mu\)，
其他周期孔由 \((C_w\setminus U)\cup H_\mu\) 上包。
因 \(x\notin\mathcal B_\mu(U)\)，同一见证严格满足
\[
 M+R_U(x)+R_{H_\mu}(x)<M+F(s)=a(\rho+s)^2.
\]
微孔 host 在 \(\operatorname{int}U_2\subset U\)，没有宏孔；
不同宏胞的外包孔分离，所以短见证在最终输入中仍严格有效。

相对原宏观参数，此改动给
\[
 \Delta\rho=2a\Lambda_0u^6+O(u^9),\qquad
 \Delta M=4a^2g\Lambda_0u^6+O(av^9).
\]
于是 \(\Delta M/(av^8)=g\Lambda_0/16+O(u)\to8/27\)；
比值 \(d/(d-M)\) 对 \(M\) 的导数为 \(1+O(a)\)，
故宏观八阶系数恰增加 \(8/27\)。
这里总质量保持为新指定 \(M\)，不是在原宏孔上直接再加微孔。

最后仍取整数 \(N=\lceil u^{-10}\rceil\)、凸外盒
\(\Omega_N=T([-N,N]^2)\)，并保留全部边界切孔。
周期积分精确给输入质量 \((2N)^2(d-M)>0\)。
宏、微见证均小于一，\(\|T^{-1}\|<2\)，所有内层孔点的指定见证保留；
边界孔总质量 \(O(aN)\)。凸盒外任意中心任意半径方体至多半体积在
分离半空间内，小 \(t\) 下不能进入严格 halo；输入密度点也进入 halo。
所以全部真实半径的比值为 \(d/(d-M)+O(a/N)\)，而非仅给指定半径的子族。
有限化误差晚于 \(av^9\)，得到二维系数
\(J_2+8/27=17638/6561\)。
\end{proof}

\begin{proof}[定理~\ref{eighth:theorem} 的证明]
命题~\ref{eighth:upper} 给全输入上极限；
命题~\ref{eighth:lower} 对每个充分小实参数给实际有限输入，
故下极限至少为 \(\widehat J_n\)。二者即得
\eqref{eighth:exact-expansion}。
\end{proof}

